\documentclass[11pt]{article}
\usepackage[a4paper,margin=25mm]{geometry}
\usepackage[T1]{fontenc}
\usepackage{lmodern,microtype,amsmath,amssymb,amsthm,mathtools}
\usepackage{xcolor,enumitem}
\usepackage[colorlinks=true,linkcolor=blue!45!black,urlcolor=blue!55!black,citecolor=blue!45!black]{hyperref}
\hypersetup{pdftitle={A converse to the Archimedes sphere theorem with only Hausdorff measure conditions},pdfauthor={Yuetong Zhang}}
\newtheorem{theorem}{Theorem}[section]
\newtheorem{lemma}[theorem]{Lemma}
\newtheorem{proposition}[theorem]{Proposition}
\newtheorem{corollary}[theorem]{Corollary}
\newcommand{\R}{\mathbb R}
\newcommand{\Sph}{\mathbb S}
\newcommand{\HH}{\mathcal H}
\newcommand{\WF}{\operatorname{WF}_{A}}
\newcommand{\WFinfty}{\operatorname{WF}}
\newcommand{\supp}{\operatorname{supp}}
\newcommand{\conv}{\operatorname{conv}}
\newcommand{\Int}{\operatorname{int}}
\newcommand{\diam}{\operatorname{diam}}
\newcommand{\Span}{\operatorname{span}}
\newcommand{\dd}{\,d}
\newcommand{\restr}{\mathbin{\vrule height 1.4ex depth 0pt width .06ex\vrule height .06ex depth 0pt width .65ex}}
\setlist{itemsep=3pt,topsep=5pt}
\setlength{\emergencystretch}{2em}
\setcounter{tocdepth}{1}
\title{A converse to the Archimedes sphere theorem\\
with only Hausdorff measure conditions}
\author{Yuetong Zhang}
\date{}
\makeatletter
\renewcommand{\@maketitle}{%
  \newpage\null\vskip 2em
  \begin{center}%
    {\LARGE\@title\par}%
    \vskip 1em
    {\large Generated by: \@author\par}%
  \end{center}%
  \vskip 1.5em
}
\makeatother
\begin{document}
\maketitle
\begin{abstract}
Let $S\subset\R^3$ be a compact connected embedded topological surface without boundary. We prove that if every admissible closed slab of a fixed width $0<h<\diam S$ contains the same finite two-dimensional Hausdorff measure $C$, then $S$ is a round sphere of radius $C/(2\pi h)$. The slab condition itself implies finite total area and places $h$ below the minimum directional width. For surfaces of finite total area and $0<h<w_*(S)$, we also prove rigidity when the slab area is constant in position for each direction and may initially depend on that direction. The main step identifies $S$ with the boundary of its convex hull. Adapting the Radon transport method of Ghomi, Howard and Lai, we establish an analytic graph criterion, regularity of transported normal offsets, and almost-everywhere first-order normals along analytic arcs. Their convex rigidity theorem then gives the sphere conclusion.
\end{abstract}

\section{Introduction}\label{sec:intro}
Archimedes' sphere theorem states that the area of a spherical zone depends only on its width: for a sphere of radius $R$, a zone of width $h$ has area $2\pi Rh$. Ghomi's fixed-width strip question \cite[Problem 4.3]{Ghomi} asks for a converse. We prove such a converse for embedded topological surfaces, with area measured by $\HH^2$.

For a compact connected surface $S\subset\R^3$ and $u\in\Sph^2$, put
\[
 a(u)=\min_S x\cdot u,\qquad b(u)=\max_S x\cdot u,\qquad
 w(u)=b(u)-a(u),\qquad w_*(S)=\min_{u\in\Sph^2}w(u).
\]
A closed slab of width $h$ in direction $u$ is \emph{admissible} if its two boundary planes meet $S$. Connectedness makes this equivalent to $a(u)\le t\le b(u)-h$ for the slab $\{t\le x\cdot u\le t+h\}$. Area means the usual normalized two-dimensional Hausdorff measure $\HH^2$.

\begin{theorem}[Fixed-width strip rigidity]\label{thm:main}
Let $S\subset\R^3$ be a compact connected embedded topological two-manifold without boundary. Suppose $0<h<\diam S$ and every admissible closed slab of width $h$ contains the same finite area $C$. Then $S$ is a round sphere of radius $C/(2\pi h)$.
\end{theorem}
\begin{proof}[Proof of Theorem~\ref{thm:main}]
Lemma~\ref{lem:originalreduction} gives finite positive total area and $h<w_*(S)$. Theorem~\ref{thm:directional} therefore gives that $S$ is a round sphere. Archimedes' strip formula gives $C=2\pi Rh$ for a sphere of radius $R$; see also \cite[Introduction]{GHL}. Hence $R=C/(2\pi h)$.
\end{proof}

Theorem~\ref{thm:main} follows from the more general directionwise statement below, together with the finite-area and minimum-width reduction in Lemma~\ref{lem:originalreduction}.
\begin{theorem}[Directionwise strip rigidity]\label{thm:directional}
Let $S\subset\R^3$ be a compact connected embedded topological two-manifold without boundary, with $0<\HH^2(S)<\infty$. Suppose $0<h<w_*(S)$ and, for each $u\in\Sph^2$, the area in every admissible closed slab of width $h$ in direction $u$ has a value $C(u)$ independent of its position. Then $S$ is a round sphere. In particular, $C(u)$ is independent of $u$.
\end{theorem}
\begin{proof}[Proof of Theorem~\ref{thm:directional}]
Proposition~\ref{prop:convexification} identifies $S$ with the boundary of the compact convex body $K$ of Lemma~\ref{lem:elementary}. Its minimum directional width exceeds $h$, and its directionwise strip property is unchanged. Theorem~\ref{thm:ghl} therefore gives that $S$ is a round sphere.
\end{proof}

The main geometric step is the convexification identity
\[
 S=\partial\conv S.
\]
Once this is established, the following convex rigidity theorem completes the proof.
\begin{theorem}[Ghomi--Howard--Lai {\cite[Theorem 1.1]{GHL}}]\label{thm:ghl}
Let $K\subset\R^3$ be a compact convex body with nonempty interior, and let $\Gamma=\partial K$. If $0<h<w_*(\Gamma)$ and the area of $\Gamma$ in admissible closed slabs of width $h$ is independent of position for each fixed direction, then $\Gamma$ is a sphere.
\end{theorem}
This theorem allows the slab constant to depend on direction and applies to arbitrary convex-body boundaries. Ghomi, Howard and Lai also prove rigidity for globally $C^2$ surfaces \cite[Theorem 1.2]{GHL}. Here the topological surface and its Hausdorff area measure provide the starting point for proving convexity and regularity.

Our proof builds on their slab periodicity and analytic Radon transport \cite[equations (3), (7) and Lemma 5.5, pp.~5, 9]{GHL}, their forced-circle construction \cite[Proposition 5.1, Steps 1--3, pp.~9--10]{GHL}, and their offset and Hopf convexification arguments \cite[Lemma 6.2 and Proposition 6.4, pp.~11--13]{GHL}. Three local results extend these mechanisms to topological surfaces: a graph criterion obtained from positive Hausdorff area and projected moments, regularity of analytic normal offsets at topological targets, and an almost-everywhere arc-normal theorem. The microlocal tools are analytic Radon recovery for compactly supported distributions \cite[Theorem 1.2]{MST} and Kashiwara's Watermelon theorem for $C^1$ supporting functions \cite[Theorem 2.1 and p.~1226]{Horm}.

Section~\ref{sec:foundation} establishes strip periodicity and analytic transport, then introduces the three local results. Section~\ref{sec:global} uses them to prove convexification and both sphere theorems. Appendices~\ref{app:graph}--\ref{app:arc} state and prove the graph, offset, and arc-normal results, respectively.

Throughout the convexification argument we work under the hypotheses of Theorem~\ref{thm:directional}; Lemma~\ref{lem:originalreduction} supplies the reduction from Theorem~\ref{thm:main}. Put $\mu=\HH^2\restr S$, $K=\conv S$, and $E=S\cap\partial K$. A point of $S$ is \emph{analytic} if $S$ agrees with a real analytic embedded surface in an ambient neighborhood of that point. An \emph{actual analytic patch} is a relatively open analytic patch of $S$. Write $B$ for the closed set of nonanalytic surface points. The strip condition is
\begin{equation}\label{strip}
 \mu\{x:t\le x\cdot u\le t+h\}=C(u)
 \quad\text{if }a(u)\le t\le b(u)-h.
\end{equation}
Analytic and ordinary wavefront sets are denoted by $\WF$ and $\WFinfty$; their covectors are always nonzero. Fourier transforms use $e^{-ix\cdot\xi}$ and inverse factor $(2\pi)^{-n}$. All measures are completed when projected analytic sets occur.

\section{Strip periodicity, analytic transport, and local regularity}\label{sec:foundation}
\subsection{Strip periodicity and microlocal tools}
\subsubsection{Full support, widths, and distributional periodicity}
\begin{lemma}\label{lem:elementary}
One has $\supp\mu=S$, and $K$ has nonempty interior. If $q=R\mu$ is the plane Radon transform, then
\begin{equation}\label{period}
 q(u,t)=q(u,t-h),\qquad a(u)+h<t<b(u),
\end{equation}
as a joint distribution on this open subset of $\Sph^2\times\R$.
\end{lemma}
\begin{proof}

Every relatively open neighborhood in $S$ contains a compact topological two-disk $D$, presented as the image of $[0,1]^2$ under a homeomorphism. Denote its opposite closed boundary edges by $L,R$ and $B,T$, and put $\delta_1=\operatorname{dist}(L,R)>0$, $\delta_2=\operatorname{dist}(B,T)>0$. With $d(x,F)=\operatorname{dist}(x,F)$, define
\[
 G(x)=\left(\frac{d(x,L)}{d(x,L)+d(x,R)},
 \frac{d(x,B)}{d(x,B)+d(x,T)}\right).
\]
Its coordinate Lipschitz constants are at most $\delta_1^{-1},\delta_2^{-1}$: for two pairs of nonnegative distances $(a,b),(a',b')$, the numerator $|ab'-a'b|$ is at most $(a'+b')|x-x'|$, and each denominator sum is at least the corresponding opposite-edge distance. The boundary of $G|_D$, in square coordinates, maps each side into itself; its straight homotopy to the identity stays on those sides. It therefore has degree one about every interior square point, so $G(D)=[0,1]^2$. The Hausdorff-measure Lipschitz inequality gives
\[
 1\le(\delta_1^{-2}+\delta_2^{-2})\HH^2(D).
\]
Thus every nonempty relatively open neighborhood has positive area, and $\supp\mu=S$.

If $K$ had empty interior, $S$ would be contained in an affine plane. Invariance of domain applied to its two-manifold charts would make $S$ open in that plane, contradicting compactness. Hence $K$ is a convex body.

The window mass
\[
 W(u,t)=\int_S\mathbf1_{\{t\le x\cdot u\le t+h\}}\,d\mu(x)
\]
is jointly Borel and bounded. In particular $C(u)=W(u,(a(u)+b(u)-h)/2)$ is measurable. For each $u$, the strip hypothesis makes $W(u,\cdot)$ constant on $a(u)<t<b(u)-h$. The closed-endpoint window agrees almost everywhere in $t$ with the corresponding convolution of the projected measure; a finite measure has at most countably many atoms. Fubini therefore makes its $t$ derivative vanish as a joint distribution on the admissible open set. That derivative is $q(u,t+h)-q(u,t)$, which gives \eqref{period}.
\end{proof}

\begin{lemma}[Finite area and strict minimum width]\label{lem:originalreduction}
Under the hypotheses of Theorem~\ref{thm:main}, one has $0<\HH^2(S)<\infty$ and $h<w_*(S)$. Thus the hypotheses of Theorem~\ref{thm:directional} hold with $C(u)=C$.
\end{lemma}
\begin{proof}
Let $d=\diam S$. A direction along a pair realizing the diameter has projection interval of length $d>h$; every other width is at most $d$. At most $\lfloor d/h\rfloor+1$ closed intervals of length $h$ contained in the diameter interval cover it. Their slabs are admissible by connectedness, so
\[
 \HH^2(S)\le (\lfloor d/h\rfloor+1)C<\infty.
\]
The topological-disk argument in Lemma~\ref{lem:elementary} shows that every nonempty relatively open surface neighborhood has positive area. In particular, $\HH^2(S)>0$.

If some width were at most $h$, continuity of $w$ on the connected direction sphere, together with $\max w=d>h$, would give a direction of width exactly $h$. Its supporting slab contains all of $S$, hence $C=\HH^2(S)$. A strictly interior slab in a direction of width greater than $h$ leaves a nonempty relatively open cap, of positive area, contradicting that equality. Therefore $h<w_*(S)$. This is the intermediate-width argument of \cite[Note 1.3]{GHL}, with full area support verified here for a topological surface.
\end{proof}

The window argument gives the joint distributional form of the periodicity in \cite[equation (3), p.~5]{GHL}. There it follows from the coarea identities \cite[equations (1)--(2), p.~4]{GHL}; here it follows directly from the Hausdorff area measure.

\subsubsection{Support conormals, normal approximation, and a set-valued criterion}
The following microlocal and variational tools relate supporting functions to regularity of projections. The supporting-function construction in Section~\ref{sec:frechet} gives their application at topological surface points, extending the use of Watermelon in \cite[Lemma 5.3 and Proposition 5.1, pp.~8--10]{GHL}.
\begin{enumerate}
\item If a distribution $f$ is supported on a closed set $F$, $p\in\supp f$, and a $C^2$ function has a nonzero differential at $p$ and a local one-sided extremum on $F$ there, its conormal and its negative belong to $\WF(f)_p$. Moreover, if $\nu$ is such an exterior conormal, Kashiwara's Watermelon theorem gives
\begin{equation}\label{watermelon}
 \xi\in\WF(f)_p\ \Longrightarrow\ 
 \xi+t\nu\in\WF(f)_p
 \quad(t\in\R,\ \xi+t\nu\ne0).
\end{equation}
H\"ormander \cite[Definition 1.1, equation (1.2), Theorem 2.1]{Horm} gives both statements, including the symmetric conormal inclusion.
The same addition law holds for $C^1$ supporting functions \cite[p.~1226, paragraph after Theorem 2.3]{Horm}. Section~\ref{sec:frechet} constructs such a function from each nonzero Fr\'echet normal in the ambient coordinates.
\item For a closed set, every nonzero proximal normal comes from a nearest-point ball touching the set from its complement. Every limiting, or Mordukhovich, normal is a limit of proximal normals at nearby points. See Rockafellar--Wets \cite[Example 6.16, Exercise 6.18(a)]{RW}. Consequently,
\begin{equation}\label{normalWF}
 N_F^{\mathrm{lim}}(p)\setminus\{0\}\subset\WF(f)_p
 \quad\text{if }\supp f=F.
\end{equation}
Indeed, the touching balls supply exterior conormals, and closedness of analytic wavefront away from zero supplies the limit.
\item Let $F:\R^n\rightrightarrows\R^m$ have locally closed graph and $(0,0)\in\operatorname{graph}F$. If its limiting normal cone at $(0,0)$ contains no nonzero covector of the form $(\xi,0)$, there are neighborhoods $V,W$ of zero and $L<\infty$ with
\begin{equation}\label{aubin}
 F(x')\cap W\subset F(x)+L|x'-x|\overline B^m
 \quad(x,x'\in V).
\end{equation}
This is the coderivative criterion for the Aubin property \cite[Theorem 9.40]{RW}; see also \cite[Definition 7, Proposition 8]{DP}.
\item Ordinary wavefront is contained in analytic wavefront. Multiplication by a smooth function does not enlarge ordinary wavefront. For a projection $\pi(x,y)=x$ proper on the distributional support, absence of all pulled-back nonzero covectors $(\xi,0)$ makes its pushforward smooth; see \cite[Corollary 6.2, Theorem 6.3]{BDH}.
\end{enumerate}

\subsubsection{Exact analytic Radon transport}
The analytic Radon correspondence of \cite[equation (7), p.~9]{GHL} takes the following form at an arbitrary nonzero frequency $\tau$. We identify $T_u^*\Sph^2$ with $u^\perp\subset\R^3$ using the Euclidean metric. For the kernel $\delta(t-x\cdot u)$, a source covector $(p,\tau u)$, $\tau\ne0$, corresponds to
\begin{equation}\label{canonical}
 (u,t=p\cdot u;\eta=-\tau(p-tu),\tau).
\end{equation}
The source is recovered uniquely by $p=tu-\eta/\tau$. The left canonical projection is an injective immersion, so the Euclidean hyperplane transform satisfies the Bolker condition. Analytic recovery under this condition \cite[Theorem 1.2]{MST}, together with the forward analytic wavefront inclusion, gives wavefront equivalence for arbitrary compactly supported distributions.

The shift $(u,t)\mapsto(u,t-h)$ leaves $(\eta,\tau)$ fixed and changes the recovered source to $p-hu$. Thus any compactly supported distribution $f$ satisfying \eqref{period} for $Rf$ obeys
\begin{equation}\label{transport}
 (p,\tau u)\in\WF(f)
 \quad\Longleftrightarrow\quad
 (p-hu,\tau u)\in\WF(f),
\end{equation}
whenever
\begin{equation}\label{admissible}
 a(u)<p\cdot u-h<p\cdot u<b(u).
\end{equation}
The functions $a,b$ specify the open data set of the Radon identity, and the equivalence holds for both signs of $\tau$. In particular, it applies to any auxiliary compactly supported distribution with support in $S$ whose Radon transform satisfies \eqref{period} on that data set. This is the translation mechanism of \cite[Lemma 5.5, p.~9]{GHL}, expressed in terms of the strict overlap condition \eqref{admissible}.

\subsection{Three local regularity theorems}
The first two theorems recover analytic surface patches from transported area singularities. The third produces first-order normals along an analytic arc in a topological surface. Their proofs occupy Appendices~\ref{app:graph}--\ref{app:arc}.
The analytic graph criterion is Theorem~\ref{thm:graph}, proved in Appendix~\ref{app:graph}.
The proof first obtains a graph from Hausdorff area and manifold topology, then proves analyticity through projected moments. This extends the regularity step in \cite[Lemma 5.4, p.~8]{GHL} to a covector plane specified in ambient coordinates.

The regularity theorem for transported analytic offsets is Theorem~\ref{thm:nofocal}, proved in Appendix~\ref{app:focal}.

The first-order arc-normal theorem is Theorem~\ref{thm:arcnormal}, proved in Appendix~\ref{app:arc}.
Here a Fr\'echet normal $m$ at $q$ means $m\cdot(x-q)\le o(|x-q|)$ on the surface as $x\to q$.

\section{Convexification and the sphere conclusion}\label{sec:global}
\begin{proposition}[Convexification]\label{prop:convexification}
Under the hypotheses of Theorem~\ref{thm:directional}, $S=\partial\conv S$.
\end{proposition}
\begin{proof}\leavevmode\par
\subsection{First-order supports and the forced circular arc}\label{sec:frechet}
\subsubsection{Supporting functions for Fr\'echet normals}
For a closed set $M$ and $p\in M$, its Fr\'echet normal cone consists of $m$ such that
\[
 m\cdot(x-p)\le o(|x-p|)\quad(x\in M,\ x\to p).
\]
Each nonzero Fr\'echet normal $m$ admits a $C^1$ supporting function. To construct it, translate $p$ to zero and set
\[
 \omega(r)=\sup\left(\{0\}\cup
 \left\{\frac{(m\cdot x)_+}{|x|}:x\in M,\ 0<|x|\le r\right\}\right).
\]
This finite nonnegative function tends to zero, including when $p$ is isolated in $M$. For a nonnegative smooth function $\varphi$ supported in $(0,\log2)$ with integral one, put
\[
 \eta(r)=\int\omega(2re^s)\varphi(s)\,ds,
 \qquad g(r)=r\eta(r).
\]
Then $\omega(r)\le\eta(r)\le\omega(4r)$ and, by integrating the derivative onto $\varphi$ in logarithmic coordinates,
$|r\eta'(r)|\le\|\varphi'\|_1\omega(4r)\to0$. Thus $g(|x|)$ is $C^1$ at zero with derivative zero, and
\[
 f(x)=m\cdot x-g(|x|)\le0\quad(x\in M),\qquad df(0)=m.
\]
H\"ormander's $C^1$ supporting-function version of Kashiwara's theorem \cite[p.~1226, paragraph after Theorem 2.3]{Horm} supplies the normal covector and the Watermelon addition rule along $m$. Realness supplies its negative. Closedness of wavefront also includes limits of Fr\'echet normals, as used in \eqref{fiberAubin}. Under a $C^1$ coordinate diffeomorphism, transpose differentiation transports these normals: the Taylor error is $o(|x-p|)$. Thus the normals obtained after analytic straightening in Theorem~\ref{thm:arcnormal} pull back to Fr\'echet normals of the original surface.

\subsubsection{A nonanalytic hull contact generates a circular arc}
We adapt the circular-arc construction of \cite[Proposition 5.1, Steps 1--3, pp.~9--10]{GHL}. Theorem~\ref{thm:arcnormal} supplies first-order normals along the arc, and Theorem~\ref{thm:nofocal} excludes the resulting focal offset.

Suppose $p\in E$ is not an analytic surface point. Its outward hull normal cone is
\[
 N=\{n:n\cdot(x-p)\le0\ \text{for all }x\in K\}.
\]
It is nonzero and pointed, since $K$ has nonempty interior (Lemma~\ref{lem:elementary}). We construct a nonconstant direction arc $u(s)$ in a fixed two-plane $P$, approaching a supporting unit normal $n$, with
\begin{equation}\label{badcircleadmissible}
 (p,u(s))\in\WF(\mu),\quad
 a(u(s))<p\cdot u(s)-h<p\cdot u(s)<b(u(s)),\quad n\cdot u(s)>0.
\end{equation}
If $N$ is a ray $\R_+n$, nonanalyticity and Theorem~\ref{thm:graph} imply a wavefront covector $\xi$ not parallel to $n$: a fiber contained in that line would avoid a transverse covector plane and make the surface analytic. Watermelon gives $\xi+\lambda n$ for large $\lambda$. Their directions approach $n$ and lie outside $N\cup(-N)$.

If $N$ contains two independent vectors, choose a linear functional $c$ strictly positive on $N\setminus\{0\}$ and the compact base $D=\{v\in N:c(v)=1\}$. Choose a proper nonempty exposed face $\{v\in D:r(v)=\beta\}$, where $r\le\beta$ on $D$, and set $\ell=\beta c-r$. Then $\ell\ge0$ on $N$, and points $n_2$ in that face and $n_1$ outside it satisfy $\ell(n_2)=0$ and $\ell(n_1)>0$. Extend these functionals from the span of $N$ to $\R^3$. Starting with the supporting wavefront covector $-n_1$ and adding $\lambda n_2$ gives an arc approaching $n=n_2/|n_2|$. Its $\ell$ value is negative, so it lies outside $N$, and pointedness keeps it outside $-N$ for large $\lambda$.

In both cases the upper strict inequality in \eqref{badcircleadmissible} follows from being outside $N$, and the lower one follows by continuity from $p\cdot n-a(n)=w(n)>h$. The direction curve is an open part of the great circle in $P$, never parallel to $n$. Transport gives the regular embedded analytic circular arc
\begin{equation}\label{forcedcircle}
 \gamma(s)=p-hu(s)\in S,\qquad (\gamma(s),u(s))\in\WF(\mu).
\end{equation}
Every point of this arc is nonanalytic. Otherwise its transported covector is a normal of an actual analytic source patch. Strict overlap persists near that point, and Theorem~\ref{thm:nofocal} makes the reverse offset through $p$ a regular analytic patch. Invariance of domain would make $p$ analytic, a contradiction.

For $q=\gamma(s)$ and any unit wavefront direction $v$ near $u(s)$, strict reverse transport puts $q+hv\in S\subset K$. Supporting by $n$ at $p$ therefore gives the cap constraint
\begin{equation}\label{angularcap}
 n\cdot v\le n\cdot u(s).
\end{equation}

\subsubsection{First-order normals fill the conormal planes}
Let $W_q=N_q^*\gamma$. A nonzero Fr\'echet normal $m$ at $q$ annihilates the tangent line of $\gamma$, so lies in $W_q$. If $m$ is independent of $u$, Watermelon addition, conicity, reality, and closedness generate $\operatorname{span}(u,m)=W_q$ inside $\WF(\mu)_q$.

If $m$ is parallel to $u$, nonanalyticity and Theorem~\ref{thm:graph} supply $\xi\in\WF(\mu)_q$ not parallel to $u$. The same operations fill $L=\operatorname{span}(u,\xi)$. Choose a unit $v\in L$ perpendicular to $u$. Both signs of small $\theta$ give wavefront directions $\cos\theta\,u+\sin\theta\,v$. Differentiating \eqref{angularcap} at its maximum $\theta=0$ gives $n\cdot v=0$. Since $n,u$ span $P$, $v$ is perpendicular to $P$, and $L=W_q$. Thus
\begin{equation}\label{normalfills}
 N_F(S,q)\ne\{0\}\quad\Longrightarrow\quad
 W_q\setminus\{0\}\subset\WF(\mu)_q.
\end{equation}

If some $q_0$ has a missing nonzero covector in $W_{q_0}$, choose an analytic normal frame along the arc and an analytic defining coordinate $T$ vanishing there with that differential at $q_0$. Realness and closedness of wavefront give a smaller subarc and neighborhood with a uniform missing double cone about $dT$. Theorem~\ref{thm:arcnormal} supplies nonzero Fr\'echet normals at almost every point of that smaller arc. A full-measure subset of the regular arc is dense in every subarc. Equation~\eqref{normalfills} fills the conormal planes on that set. Approximate each prescribed nonzero covector at $q_0$ along the continuous conormal frame; closedness of wavefront fills the plane at $q_0$, contradicting the missing covector. We conclude that the full conormal plane is present at every point of a smaller circular arc.

\subsubsection{Conormal transport and focal exclusion}
Transport of the filled conormal planes produces the two-parameter family of \cite[Proposition 5.1, Step 3, p.~10]{GHL}. Its reverse normal offset has rank one, which yields the contradiction through Theorem~\ref{thm:nofocal}.
Choose a unit $e\perp P$ and set
\[
 v(s,\theta)=\cos\theta\,u(s)+\sin\theta\,e,
 \qquad Q(s,\theta)=\gamma(s)+hv(s,\theta).
\]
For a short $s$ interval and sufficiently small $\theta$, the conormal condition and reverse strict transport put $Q$ in $S$. Because $u'=-\gamma'/h$,
\[
 Q_s=(1-\cos\theta)\gamma',\qquad
 Q_\theta=h(-\sin\theta\,u+\cos\theta\,e).
\]
These vectors are independent for small nonzero $\theta$. Restrict to an embedded analytic patch. Invariance of domain makes it an actual relatively open patch of $S$. Its normal $-v$ has the strictly admissible offset
\[
 Q+h(-v)=\gamma(s),
\]
of rank one. This contradicts Theorem~\ref{thm:nofocal}. Therefore every hull contact is analytic.

\subsection{Convexity from analytic hull contacts}
We now apply the local convexification argument of \cite[Proposition 6.4, pp.~12--13]{GHL} on the analytic neighborhood of the contact set. Topological separation fixes the outward normal, the offset identity gives a mean-curvature inequality at noncontacts, and Hopf's lemma shows that the contact set is all of $S$.
The topological separation theorem gives exactly two components $G,O$ of $\R^3\setminus S$, with $\partial G=\partial O=S$; $G$ is bounded and $O$ unbounded. The separation theorem applies to embedded topological surfaces \cite[Corollary 3 and Theorem 4]{Lemmens}.

At an actual analytic graph patch, its two local sides belong one each to $G$ and $O$, since both components approach every point. One has $G\subset\Int K$: a point outside $K$ has a separating outward ray disjoint from $K$, hence from $S$, connecting it to infinity, and $G$ is open.

The compact contact set $E$ is separated from the closed nonanalytic locus, so it has an analytic neighborhood in $S$. Let $n_{\rm out}$ denote the analytic normal pointing from $G$ into $O$ there. At contacts it agrees with the hull-support normal, since an outward supporting ray lies in $O$. Put $\delta=\min w-h>0$. The continuous support defect
\[
 d(x)=b(n_{\rm out}(x))-x\cdot n_{\rm out}(x)
\]
vanishes on $E$ and is positive at nearby noncontacts. Shrink the neighborhood so $d<\delta/2$. For a noncontact its inward offset $x-hn_{\rm out}$ is strictly admissible, since
\[
 x\cdot n_{\rm out}-a(n_{\rm out})=w(n_{\rm out})-d(x)>h,
 \qquad x\cdot n_{\rm out}<b(n_{\rm out}).
\]
Apply Theorem~\ref{thm:nofocal} with source normal $-n_{\rm out}$. If $\kappa_1,\kappa_2$ are the outward principal curvatures, positive on a round sphere, its amplitude identities give
\[
 1-h\kappa_i>0,
 \qquad (1-h\kappa_1)(1-h\kappa_2)=1.
\]
Arithmetic--geometric mean gives $H_{\rm out}=\kappa_1+\kappa_2\le0$ at every such noncontact, as in \cite[Lemma 6.2, p.~12]{GHL}. For the inward source normal $-n_{\rm out}$, our $A=-dn$ equals $dn_{\rm out}$, the shape operator used there; our $H_{\rm out}$ is twice their mean curvature $H$.

If $E\ne S$, connectedness supplies a boundary point of $E$ approached by noncontacts. Choose an analytic graph chart near it, with vertical coordinate oriented so that $n_{\rm out}$ is the upward normal throughout the chart. In the two-dimensional base plane, choose the projection of a noncontact point sufficiently close to the projected contact set. The closed base disk centered there, with radius equal to its distance to that set, stays in the chart and in the analytic neighborhood. Its open disk consists of projections of noncontacts, and its boundary contains the projection $y_0$ of a contact. Write the surface as $y=f(x)$ and its supporting affine function at that contact as $\ell(x)$. The graph is below $\ell$; equality would itself be a hull contact. Thus $w=f-\ell<0$ inside the base disk, while $w(y_0)=0$ and $Dw(y_0)=0$. With the chosen upward outward normal, the mean-curvature inequality gives
\[
 \left(\frac{\delta_{ij}}{\sqrt{1+|Df|^2}}
 -\frac{f_if_j}{(1+|Df|^2)^{3/2}}\right)w_{ij}\ge0.
\]
This operator is uniformly elliptic on the small closed disk. Hopf's boundary point lemma \cite[Chapter 3]{GT} gives a positive outward normal derivative of $w$ at $y_0$, contrary to $Dw(y_0)=0$. Hence $S=E$. The inclusion into the connected convex boundary is open by invariance of domain and closed by compactness, so $S=\partial K$. It is now a real-analytic convex surface.

This proves Proposition~\ref{prop:convexification}.
\end{proof}
\appendix
\section{Proof of the analytic graph criterion}\label{app:graph}
\begin{theorem}[Analytic graph criterion]\label{thm:graph}
Let $S$ be a closed embedded topological two-manifold without boundary in $\R^3$, with locally finite $\HH^2$ measure, and put $\mu=\HH^2\restr S$. In coordinates $(x,y)\in\R^2\times\R$, let $p=(0,0)\in S$. Suppose
\begin{equation}\label{horizontal}
 (p;\xi,0)\notin\WF(\mu)\qquad(0\ne\xi\in\R^2).
\end{equation}
Then $S$ is a real analytic graph $y=f(x)$ near $p$.
\end{theorem}
\begin{proof}[Proof of Theorem~\ref{thm:graph}]
Choose an ambient neighborhood of $p$ in which $S$ is relatively closed, and regard $\mu$ as a distribution there. By Lemma~\ref{lem:elementary}, every topological manifold neighborhood has positive area, so the relative support of $\mu$ is $S$. Closedness of the analytic wavefront set and compactness of the horizontal unit covectors allow us to shrink to an ambient neighborhood $U$ of $p$ in which every nonzero horizontal covector is absent from the analytic wavefront set. The ordinary wavefront set is contained in the analytic wavefront set. Thus $S\cap U$ is relatively closed, $\supp_U\mu=S\cap U$, and every nonzero horizontal covector is also absent from the ordinary wavefront set. The ordinary pushforward theorem then gives a smooth projected density for every $\chi\in C_c^\infty(U)$ under $\pi(x,y)=x$.

\emph{Projected densities are at least one at projected support points.}
Fix $q\in S\cap U$, choose a nonnegative cutoff $\chi$ equal to one on a neighborhood $O$ of $q$, and write $\pi_*(\chi\mu)=g(x)\dd x$. To prove $g(\pi q)\ge1$, suppose that a base ball $B$ around $\pi q$ satisfies $g\le c<1$. Set $E=S\cap O\cap\pi^{-1}(B)$ and $P=\pi(E)$. For every Borel $D\subset B$, the projection inequality gives
\[
 \int_Dg\dd x\ \ge\ \mu(E\cap\pi^{-1}(D))
 \ \ge\ \HH^2(P\cap D).
\]
The projected set $P$ is Lebesgue measurable, and the inequality implies $g\ge1$ almost everywhere on $P$. Consequently $|P|=0$. Absolute continuity and positivity of the pushforward, taken in the completed measures, give $\mu(E)=0$. Since $E$ is a relative neighborhood of $q$, full support gives $\mu(E)>0$. This contradiction proves $g(\pi q)\ge1$.

\emph{Projection is open and its local fibers are finite.}
For any small relative neighborhood $A$ of $q$, choose the preceding cutoff so that $S\cap\supp\chi\subset A$. The bound $g(\pi q)\ge1$ and continuity give $g>0$ on a base neighborhood of $\pi q$. Positivity and compact support imply
\[
 \supp\pi_*(\chi\mu)=\pi(\supp\chi\mu)\subset\pi(A),
\]
so $\pi|_{S\cap U}$ is locally open.

Next fix $\chi_0\ge0$ supported in $U$ and equal to one on a smaller neighborhood $U_0$, and let $g_0$ be its projected density. Given $N$ distinct points of $S\cap U_0$ in a common fiber over $x_0$, choose disjoint cutoffs $0\le\chi_j\le1$ supported in $U_0$ and equal to one near those points. Their smooth densities satisfy
\[
 N\le\sum_{j=1}^N g_j(x_0)\le g_0(x_0).
\]
The first inequality is the projected-density bound. The second follows from $\sum_j\chi_j\le\chi_0$ and continuity of the densities. Hence every such fiber is finite, and $p$ is isolated in its local vertical fiber.

\emph{A continuous selection gives the entire local surface.}
Choose $b>0$ and then a small base ball $B_a$ so that the cylinder lies in $U_0$ and
\[
 S\cap(\{0\}\times[-2b,2b])=\{p\},\qquad
 S\cap\bigl(B_a\times\{y:b\le|y|\le2b\}\bigr)=\varnothing.
\]
The first condition follows from isolation in the vertical fiber; local closedness and compactness give the second. By openness of projection,
\[
 F(x)=\{y\in[-2b,2b]:(x,y)\in S\}
\]
is nonempty for every $x$ in a smaller $B_a$. Each $F(x)$ is a compact subset of $(-b,b)$. Define $f(x)=\min F(x)$. Closedness of the graph gives lower semicontinuity of $f$, since a subsequential limit of minimizing values belongs to $F(x)$. Openness of projection at $(x,f(x))$ gives upper semicontinuity, since nearby fibers meet every fixed small interval about $f(x)$. Thus $f$ is continuous. Invariance of domain applied to $x\mapsto(x,f(x))$ makes its image relatively open in $S$. Since this image contains $p$, the surface agrees with the continuous graph in an ambient neighborhood of $p$.

\emph{Analytic projected densities give an analytic graph.}
Shrink to an open graph cylinder $V=B\times(-b,b)$ such that the graph over compact subsets of $B$ stays away from its top and bottom. Restriction to $V$ preserves the analytic wavefront condition, and projection $V\to B$ is proper on the support of the restricted measure. Define the projected densities by
\[
 q_0\,dx=\pi_*(\mu\restr V),\qquad
 q_1\,dx=\pi_*(y\mu\restr V).
\]
The analytic proper-pushforward theorem \cite[Corollary 5.4.10 and Remark 5.4.11]{Schultka} makes $q_0,q_1$ real analytic on $B$. Support on the continuous graph gives the measure identity $q_1\,dx=fq_0\,dx$. Continuity makes this identity pointwise, and the projected-density bound gives $q_0\ge1$. Hence $f=q_1/q_0$ is real analytic.
\end{proof}

\paragraph{Remark.}
The hypotheses enter at successive stages of the proof. Wavefront avoidance gives smooth projected densities. Unweighted area bounds these densities below by one at projected support points, yielding openness of projection and a bound on the number of local sheets. Manifold topology then makes the continuous minimum selection a neighborhood of the surface. Finally, analytic wavefront avoidance makes the projected moments analytic and hence gives an analytic graph.

\section{Positive area and regular offsets}\label{app:focal}\label{sec:focal}
\subsection{Normal-offset amplitudes}
Let $T_hg(u,t)=g(u,t-h)$ and $U_h=e^{-ih|D|}$. The positive-frequency Radon inversion operator is
\[
 B_+g(x)=(2\pi)^{-3}\int_{\Sph^2}\int_0^\infty
 e^{i\lambda x\cdot u}\lambda^2\widehat g(u,\lambda)
 \,d\lambda\,d\sigma(u).
\]
The Fourier slice identity gives $B_+R=\mathrm{Id}$ and $B_+T_hR=U_h$. A high-frequency cutoff changes the result by a smooth term. Since the positive-frequency canonical relation is a graph, localization of \eqref{period} gives
\begin{equation}\label{exceptionhalf}
 \mu-U_h\mu\text{ is microlocally smooth at }(x,\lambda u),\quad\lambda>0,
\end{equation}
under \eqref{admissible}. The unique inverse half-wave source at $(x,\lambda u)$ is $x-hu$. This identifies the source patch and permits its localization by a fixed smooth cutoff, with multiplier and phase determined by $U_h$.

For the oscillatory tests used below, let $\psi_\omega$ have unit gradient and let $\omega$ range over a compact angular subset. Shrink the target neighborhood so that
$\{x-h\nabla\psi_\omega(x):x\in\supp\chi,\ \omega\text{ in that subset}\}$
is compactly contained in the analytic source neighborhood, and choose a source cutoff equal to one near this set. The inverse canonical graph separates the remaining source microlocally, while directions outside the source conormal are smooth. These fixed cutoffs and the compact angular restriction make strict admissibility uniform in the large-frequency limit.

The following frequency localization applies to each oscillatory test below. After writing $\xi=\lambda\eta$, the half-wave phase is $\Phi=(x-y)\cdot\eta-h|\eta|-\psi_\omega(x)$. Choose a smooth radial cutoff supported in a compact annulus about $|\eta|=1$ and equal to one on a smaller annulus. Then
\[
 d_x\Phi=\eta-d\psi_\omega(x),\qquad |d\psi_\omega(x)|=1.
\]
On the complementary frequency region, $|d_x\Phi|$ is bounded away from zero and grows proportionally to $|\eta|$ for large $|\eta|$. Repeated integration by parts in $x$, using
\[
 L_x=\frac{(\eta-\nabla\psi_\omega)\cdot\nabla_x}
 {i\lambda|\eta-\nabla\psi_\omega|^2},\qquad
 L_xe^{i\lambda\Phi}=e^{i\lambda\Phi},
\]
gives an $O(\lambda^{-N})$ bound for every $N$, after the $(\lambda/(2\pi))^3$ prefactor, uniformly on the compact angular set. The large-frequency tails are integrable after sufficiently many integrations. A fixed low-frequency cutoff in $\xi$ produces a smooth term and satisfies the same estimate against the tests. Stationary phase therefore applies on the compact annulus. Analytic Radon transport identifies the source of the area measure first, and these smooth cutoffs evaluate its half-wave contribution.

For an actual analytic source patch $Q$ with analytic unit normal $n$, the strict source-overlap condition is
\begin{equation}\label{strictsourceoffset}
 a(n(y))<y\cdot n(y)<y\cdot n(y)+h<b(n(y))\quad(y\in Q).
\end{equation}
Normal-wavefront transport places $F(y)=y+hn(y)$ in $S$. At regular points of $F$, invariance of domain identifies its image as a relatively open analytic patch of $S$.
\subsubsection{The amplitude at a regular offset}
The local half-wave amplitude of unweighted area determines the sign and Jacobian of a regular offset. For nonzero Gaussian curvature, the corresponding principal-factor and Jacobian identities occur in the proof of \cite[Lemma 6.2, pp.~11--12]{GHL}, through paired Morse singularities, equal Gaussian-curvature magnitudes, and preservation of the Morse index. The calculation below also covers a vanishing principal curvature. Section~\ref{sec:nofocal} treats focal offsets by positivity of the target measure.
\begin{lemma}[Positive area fixes both phase and offset Jacobian]\label{lem:offsetamplitude}
Suppose $\mu$ equals unweighted area in ambient neighborhoods of analytic patches $Q$ and $P$. Choose an analytic unit normal $n$ on $Q$, and suppose
\[
 F(y)=y+hn(y),\qquad A_Q=-dn,
 \qquad J=\det(I-hA_Q)
\]
is a local diffeomorphism onto $P$. Assume \eqref{exceptionhalf} holds at $(F(y),\lambda n(y))$ for $\lambda>0$. If $\kappa_1,\kappa_2$ are the eigenvalues of $A_Q$, then
\begin{equation}\label{positiveoffset}
 1-h\kappa_i>0\ (i=1,2),\qquad
 (1-h\kappa_1)(1-h\kappa_2)=1.
\end{equation}
The conclusion also holds when a principal curvature is zero.
\end{lemma}
\begin{proof}
Set $j_i=1-h\kappa_i$. Choose the target normal by $n_P(F(y))=n(y)$, and let $\psi$ be signed distance to $P$ in this normal direction. Thus $\psi=0$ on $P$, $d\psi=n_P$ there, and $|d\psi|=1$. Let $\chi\ge0$ be any smooth test with sufficiently small compact ambient support near the target. The exact local area and the positive-cone regularity give, for a source cutoff equal to one at the corresponding points,
\begin{equation}\label{offsettest}
 \int_P\chi\dd\HH^2
 =\langle e^{-ih|D|}\mu_Q,\chi e^{-i\lambda\psi}\rangle
   +O(\lambda^{-\infty}),\qquad\lambda\to+\infty.
\end{equation}
The support of $\chi$ is chosen so that its positive $d\psi$ graph lies entirely in the open regularity cone, giving the rapidly decaying error above.

With $\xi=\lambda\eta$, the phase on the right is
\[
 \Phi(x,y,\eta)=(x-y)\cdot\eta-h|\eta|-\psi(x),
\]
with prefactor $(\lambda/(2\pi))^3$. By the frequency localization above, the compact annulus captures the contribution up to $O(\lambda^{-\infty})$. Its critical manifold has dimension two and is given by $x=F(y)$, $\eta=n(y)$, with critical value zero. Keeping the source coordinates free, we apply stationary phase in the six transverse variables. The normal position-frequency block is
\[
 \begin{pmatrix}0&1\\1&0\end{pmatrix}.
\]
In common orthonormal principal directions, $A_P=A_Q(I-hA_Q)^{-1}$, and the signed-distance Hessian has eigenvalues $-\kappa_i/j_i$. Consequently the tangential position-frequency blocks are
\[
 H_i=\begin{pmatrix}\kappa_i/j_i&1\\1&-h\end{pmatrix},
 \qquad \det H_i=-1/j_i.
\]
Each block has signature zero when $j_i>0$ and signature $-2$ when $j_i<0$. Thus, if $m$ is the number of negative $j_i$, the transverse Hessian has signature $-2m$ and absolute determinant $1/|J|$.

Stationary phase \cite[Section 7.7]{HormSP} therefore gives the principal term
\[
 \int_Q\chi(F(y))e^{-i\pi m(y)/2}\sqrt{|J(y)|}\dd\HH^2(y)
\]
in \eqref{offsettest}; the next term is $O(\lambda^{-1})$. Since $\dd\HH^2(F(y))=|J(y)|\dd\HH^2(y)$, varying $\chi$ and taking $\lambda\to\infty$ yields
\[
 1=\frac{e^{-i\pi m/2}}{\sqrt{|J|}}.
\]
The positive real coefficient forces $m=0$ among the possibilities $m=0,1,2$. Hence both $j_i$ are positive and $J=1$, proving \eqref{positiveoffset}. The Hessian blocks remain nondegenerate when $\kappa_i=0$, so the argument includes zero principal curvatures.
\end{proof}

\subsection{Regularity of analytic offsets}\label{sec:nofocal}
\begin{theorem}[Regularity of transported analytic offsets]\label{thm:nofocal}
Under the standing hypotheses of Theorem~\ref{thm:directional}, let $Q$ be a connected actual analytic patch of $S$ with chosen analytic unit normal $n$, and assume $a(n(y))<y\cdot n(y)<y\cdot n(y)+h<b(n(y))$ throughout $Q$. By Lemma~\ref{lem:elementary} and analytic Radon transport, $\mu=\HH^2\restr S$ satisfies the local half-wave comparison \eqref{exceptionhalf} in the strict overlap region. Put $F(y)=y+hn(y)$ and $A_Q=-dn$, and let $\kappa_1,\kappa_2$ be the eigenvalues of $A_Q$. Then $F$ is regular everywhere, its image is locally an actual analytic patch of $S$, and
\[
 1-h\kappa_i>0\ (i=1,2),\qquad
 \det(I-hA_Q)=1.
\]
\end{theorem}
\begin{proof}
To prove Theorem~\ref{thm:nofocal}, we use positivity of the target area measure to rule out the two possible focal ranks. The argument applies directly to rough targets. For comparison, \cite[Lemma 6.3, p.~12]{GHL} treats degenerate curvature by developable straight rulings and global $C^2$ continuation.

\subsubsection{Growth and homogeneous limits from scalar marginals}
\begin{lemma}[Positive Fourier-cone compactness]\label{lem:positivecone}
Let $\nu$ be a positive compactly supported finite measure on $\R^m$, and let $0<\alpha\le1$. Suppose on a nonempty open cone $U$, uniformly on compact angular subsets,
\begin{equation}\label{positiveconeasymptotic}
 \widehat\nu(\lambda u)
 =e^{-i\pi\alpha/2}A(u)\lambda^{-\alpha}
  +O(\lambda^{-\alpha-1}),\qquad \lambda\to+\infty,
\end{equation}
where $A$ is smooth and positive. Then there is a nonzero positive $\alpha$-homogeneous Radon measure $\rho$ on $\R^m$ such that
\[
 \rho(tD)=t^\alpha\rho(D),\qquad
 \widehat\rho(\xi)=e^{-i\pi\alpha/2}
 A(\xi/|\xi|)|\xi|^{-\alpha}\quad\hbox{on }U.
\]
The equality on $U$ is distributional and is an equality of smooth functions there.
\end{lemma}
\begin{proof}
Fix a unit $u_0\in U$ and the positive scalar marginal $\sigma=(x\mapsto x\cdot u_0)_*\nu$. Reality gives its negative-frequency expansion. With a compact smooth cutoff $\eta=1$ near zero,
\[
 \mathcal F[\eta(t)t_+^{\alpha-1}](\lambda)
 =\Gamma(\alpha)e^{-i\pi\alpha/2}\lambda^{-\alpha}
   +O(\lambda^{-N})\quad(\lambda>0)
\]
for every $N$; at $\alpha=1$ this is the cutoff Heaviside formula with leading term $1/(i\lambda)$. After subtracting $A(u_0)\eta(t)t_+^{\alpha-1}/\Gamma(\alpha)$ from $\sigma$, the Fourier remainder is integrable and its inverse is bounded and continuous. Thus $\sigma([-r,r])=O(r^\alpha)$ near zero. Positivity and $B_r(0)\subset\{|x\cdot u_0|\le r\}$ give
\begin{equation}\label{conegrowth}
 \nu(B_r(0))\le Cr^\alpha\qquad(r>0),
\end{equation}
where finiteness extends the bound to all radii.

The positive rescalings $\nu_\varepsilon(D)=\varepsilon^{-\alpha}\nu(\varepsilon D)$ satisfy the same bound on every centered ball $B_R(0)$, uniformly in $\varepsilon$. This gives vague compactness, and a subsequence converges to a positive Radon measure $\nu_\infty$. Applying \eqref{conegrowth} on dyadic annuli controls Schwartz tails uniformly, so the convergence also holds in tempered distributions. Since
\[
 \widehat\nu_\varepsilon(\xi)
 =\varepsilon^{-\alpha}\widehat\nu(\xi/\varepsilon),
\]
the locally uniform cone expansion yields the stated nonzero Fourier coefficient for $\nu_\infty$ on $U$.

Define normalized dilations $D_t\zeta(D)=t^{-\alpha}\zeta(tD)$ and average them logarithmically:
\[
 \rho_T=\frac1{2T}\int_{-T}^{T}D_{e^s}\nu_\infty\dd s.
\]
Every integrand satisfies the same ball-growth bound and has the same Fourier restriction on $U$, because the coefficient is homogeneous in frequency. A subsequence of $\rho_T$ therefore converges vaguely and in tempered distributions to a positive measure $\rho$ with that Fourier coefficient. The growth bound shows that shifting the averaging interval by a fixed $a$ changes any Schwartz pairing by $O(|a|/T)$. Hence $D_{e^a}\rho=\rho$ for every $a$, giving the required homogeneity. Its nonzero cone coefficient ensures that $\rho\ne0$.
\end{proof}

\begin{lemma}[Bounded complex continuation from positivity]\label{lem:complexcone}
Let $m\ge2$ and $\rho$ be a positive $\alpha$-homogeneous Radon measure.
\begin{enumerate}
\item If $0<\alpha<1$ and $\widehat\rho=e^{-i\pi\alpha/2}g$ distributionally on an open cone $U$, where $g\in C^\infty(U)$ is real and strictly positive, then for every unit $\xi_0\in U$ and perpendicular unit $e$, the function $g(\xi_0+ze)$, initially defined for real $z$ near zero, extends holomorphically and remains bounded on some vertical strip $|\operatorname{Re}z|<\delta$, including its entire imaginary axis.
\item If $\alpha=1$, the Fourier value $\widehat\rho(\xi)=-iA/|\xi|$, with $A>0$, is impossible on any nonempty open cone.
\end{enumerate}
\end{lemma}
\begin{proof}
Homogeneity gives the polar decomposition $\dd\rho(rv)=r^{\alpha-1}\dd r\dd\beta(v)$, where $\beta$ is a finite positive measure on $\Sph^{m-1}$, and implies $\rho(\{0\})=0$. For $0<\alpha<1$, radial exponential regularization gives
\[
 \widehat\rho(\xi)=\Gamma(\alpha)\int
 e^{-i\pi\alpha\operatorname{sgn}(\xi\cdot v)/2}
 |\xi\cdot v|^{-\alpha}\dd\beta(v).
\]
The formula holds in distributions, with fractional-power kernels locally integrable in $\xi$. After multiplication by $e^{i\pi\alpha/2}$, its imaginary part is the positive distribution
\[
 \Gamma(\alpha)\sin(\pi\alpha)
 \int_{\xi\cdot v<0}|\xi\cdot v|^{-\alpha}\dd\beta(v).
\]
Its vanishing on $U$ places $\supp\beta$ in the closed dual cone: $\xi\cdot v\ge0$ for every $\xi\in U$. Indeed, if $\xi\cdot v<0$ at a pair in $U\times\supp\beta$, the same inequality holds in neighborhoods of that pair, where a nonnegative test gives a strictly positive integral. Choosing $B_{2\delta}(\xi_0)\subset U$ therefore gives $\xi_0\cdot v\ge2\delta$ on the angular support. Consequently
\[
 g(\xi_0+ze)=\Gamma(\alpha)\int
 ((\xi_0+ze)\cdot v)^{-\alpha}\dd\beta(v)
\]
is holomorphic and bounded for $|\operatorname{Re}z|<\delta$, with the power branch defined on the right half-plane. For real $z$ near zero, the integral represents the Fourier distribution and agrees pointwise with its prescribed smooth representative $g$. It therefore gives the asserted extension.

At $\alpha=1$, the radial Fourier boundary-value formula is
\[
 \widehat\rho(\xi)=\int\left[
 \pi\delta(\xi\cdot v)-i\operatorname{PV}\frac1{\xi\cdot v}
 \right]\dd\beta(v).
\]
The real part is a positive mixture of hyperplane measures. Its vanishing on $U$ places each hyperplane $v^\perp$, $v\in\supp\beta$, outside $U$: a positive test in a small ball meeting such a hyperplane has positive hyperplane integral, continuously for nearby $v$. An interior ball $B_{2\delta}(\xi_0)\subset U$ then gives $|\xi_0\cdot v|\ge2\delta$ on the angular support. This bound applies to both signs of $\xi_0\cdot v$. The function
\[
 F(z)=\int\frac1{(\xi_0+ze)\cdot v}\dd\beta(v)
\]
is therefore holomorphic and bounded for $|\operatorname{Re}z|<\delta$. The assumed cone coefficient gives $F(z)=A/\sqrt{1+z^2}$ for small real $z$. Continuing the square-root branch that equals one at $z=0$ along $z=it$, $0\le t<1$, makes the right side diverge as $t\uparrow1$. The boundedness of $F$ gives a contradiction.
\end{proof}

\subsubsection{Rank-zero focal patches}
At regular points of the offset $F$ in Theorem~\ref{thm:nofocal}, invariance of domain and Lemma~\ref{lem:offsetamplitude} give $J=1$ and positive principal factors. The continuous function $J$ therefore takes values in $\{0,1\}$ on the connected patch $Q$ and is constant. To prove $J=1$, suppose that $J=0$. If $dF$ has rank zero everywhere, then $F\equiv p$ and $Q$ is an open subset of the sphere of radius $h$ centered at $p$, with inward normal $n(y)=(p-y)/h$.

Translate $p$ to zero. On a smaller open cone of normals, choose a smooth source cutoff equal to one near the corresponding points $-hu$ and supported away from the opposite critical points $hu$. Let $f$ be the localized sphere area measure. Its local lower-height density is $2\pi h$, so stationary phase on the sphere gives, uniformly on that cone,
\[
 \widehat f(\xi)=-\frac{2\pi ih}{|\xi|}e^{ih|\xi|}
    +O(|\xi|^{-\infty}).
\]
Let $\chi\ge0$ be a fixed smooth cutoff near zero with $\chi(0)>0$. Strict transport and uniqueness of the half-wave source imply that $\chi(\mu-e^{-ih|D|}f)$ has a rapidly decreasing Fourier transform on a smaller cone. The half-wave multiplier cancels the exponential above. Multiplication by $\chi$ becomes Fourier convolution, whose tails outside the cone decay rapidly. Expanding $|\lambda u-\eta|^{-1}$ in the remaining integral gives
\begin{equation}\label{sphericalfocalcoefficient}
 \widehat{\chi\mu}(\lambda u)
 =-2\pi ih\chi(0)\lambda^{-1}+O(\lambda^{-2}).
\end{equation}
Here the normalization is $(2\pi)^{-3}\int\widehat\chi=\chi(0)$, and $\chi\mu$ is a fixed positive measure in $\R^3$. Lemma~\ref{lem:positivecone} with $\alpha=1$, followed by the endpoint assertion of Lemma~\ref{lem:complexcone}, contradicts \eqref{sphericalfocalcoefficient}. This rules out rank zero.

\subsubsection{The Fourier coefficient at an analytic focal curve}
If $dF$ has rank one at some point, shrink to a constant-rank patch. Its image is a regular analytic curve $\gamma(s)$ parametrized by arclength. Write $T=\gamma'$ and choose an analytic Bishop normal frame $e_1,e_2$ satisfying
\[
 e_i'=-k_iT,\qquad \gamma''=k_1e_1+k_2e_2,
 \qquad \mathbf k=(k_1,k_2).
\]
Since $n\cdot dF=0$, the source normal is perpendicular to $T$. Along a fiber, the relation $y=\gamma-hn$ shows that the normal varies nontrivially and provides an angular coordinate. Thus
\[
 u(s,\theta)=\cos\theta\,e_1(s)+\sin\theta\,e_2(s),\qquad
 y(s,\theta)=\gamma(s)-hu(s,\theta),
\]
\[
 D(s,\theta)=1+h\mathbf k(s)\cdot(\cos\theta,\sin\theta),
 \qquad y_s=DT,\quad y_\theta=-h\partial_\theta u.
\]
Regularity of the source gives $D\ne0$ and area element $h|D|\dd s\dd\theta$. Shrink the parameter rectangle so that $D$ has a fixed sign.

In a small tubular neighborhood of the target curve, use the fixed normal-coordinate map
\[
 x=\gamma(s)+z_1e_1(s)+z_2e_2(s),\qquad\Psi(x)=(z_1,z_2).
\]
For a planar unit vector $\omega=(\cos\theta_0,\sin\theta_0)$, set $\psi_\omega=\omega\cdot\Psi$. The Bishop equations give $\nabla\psi_\omega=u(s,\theta_0)$ and $|\nabla\psi_\omega|=1$ throughout the neighborhood. On the curve, the only nonzero Hessian entry is
\[
 \nabla^2\psi_\omega(T,T)=-\mathbf k(s)\cdot\omega.
\]
Choose a nonnegative smooth compactly supported target cutoff $\chi$ with $\int\chi(\gamma(s))\dd s>0$. Restrict to a smaller angular interval so that the positive $d\psi_\omega$ graph over its support lies entirely in the microlocal regularity region. Uniqueness of the inverse half-wave source supplies a fixed source cutoff for this interval. Consequently
\[
 \nu=\Psi_*(\chi\mu)
\]
is a fixed positive planar measure with the following directional Fourier expansions.

The half-wave comparison has phase
\[
 (x-y(s,\theta))\cdot\eta-h|\eta|-\psi_\omega(x),
\]
with prefactor $(\lambda/(2\pi))^3$ after frequency scaling. As in the regular-offset calculation of Lemma~\ref{lem:offsetamplitude}, $d_x\Phi=\eta-d\psi_\omega$ localizes the integral to a fixed compact annulus about $|\eta|=1$, with an $O(\lambda^{-\infty})$ complementary contribution uniformly in $\omega$. Its critical manifold is given by $x=\gamma(s)$, $\theta=\theta_0$, $\eta=u(s,\theta_0)$, with free coordinate $s$. In the frame $(u,\partial_\theta u,T)$, the Hessian in the seven transverse variables has blocks
\[
 \begin{pmatrix}0&1\\1&0\end{pmatrix},\qquad
 \begin{pmatrix}0&1&0\\1&-h&h\\0&h&-h\end{pmatrix},\qquad
 \begin{pmatrix}\mathbf k\cdot\omega&1\\1&-h\end{pmatrix}.
\]
The middle block has determinant $h$ and signature $-1$. The last has determinant $-D$, with signature zero for $D>0$ and $-2$ for $D<0$. Thus the full transverse Hessian has absolute determinant $h|D|$ and signature $-1$ or $-3$. Including the source area element, stationary phase gives
\begin{equation}\label{curvedfocalcoefficient}
 \begin{split}
 \widehat\nu(\lambda\omega)
 ={}&\sqrt{2\pi h}\,\lambda^{-1/2}
 \int\chi(\gamma(s))\sqrt{|1+h\mathbf k(s)\cdot\omega|}\,
 e^{-i\pi(1+2\mathbf1_{D<0})/4}\dd s\\
 &{}+O(\lambda^{-3/2}).
 \end{split}
\end{equation}
Parameter-dependent stationary phase on this compact frequency region gives a uniform remainder, together with any fixed number of angular derivatives, on compact angular subintervals. The target and source cutoffs are fixed throughout the limit $\lambda\to\infty$. Away from a fixed neighborhood of the compact stationary set, integration by parts gives rapid decay. Together with the annular tail estimate, this proves the stated expansion for the full integral.

\subsubsection{Exclusion of rank-one focal patches}
The sign of $D$ determines the contradiction with positivity. For $D<0$, a scalar marginal of \eqref{curvedfocalcoefficient} has phase $e^{-3\pi i/4}$ and positive coefficient. Its negative-frequency expansion follows from reality. Since $\mathcal F[(-t)_+^{-1/2}](\lambda)=\sqrt\pi e^{i\pi/4}\lambda^{-1/2}$ for $\lambda>0$, Fourier inversion gives a density equal to a negative multiple of $(-t)_+^{-1/2}$ plus a bounded continuous remainder, contradicting positivity.

For $D>0$, set $w(s)=\chi(\gamma(s))$. Lemma~\ref{lem:positivecone} with $\alpha=1/2$ gives a positive homogeneous planar measure whose extreme-phase coefficient is
\begin{equation}\label{curvedhomogeneous}
 G(\xi)=\sqrt{2\pi h}\int w(s)
 \frac{\sqrt{|\xi|+h\mathbf k(s)\cdot\xi}}{|\xi|}\dd s.
\end{equation}
Lemma~\ref{lem:complexcone} gives a bounded holomorphic continuation along $\xi=\omega_0+itv_0$, $0\le t\le1$, where $\omega_0$ is an interior unit cone direction and $v_0$ is a perpendicular unit vector.

Suppose the curvature is nonzero at a point $s_0$ of the source interval. Choose $\omega_0$ in the angular interval so that $\mathbf k(s_0)\cdot v_0\ne0$, and choose the target cutoff at the outset with sufficiently small support that $\mathbf k(s)\cdot v_0$ has a fixed nonzero sign on the support of $w$. This preserves $D>0$. Along the complex path, set
\[
 z(t)=\sqrt{1-t^2},\qquad
 b_s(t)=z(t)+h\mathbf k(s)\cdot\omega_0
                +ith\mathbf k(s)\cdot v_0.
\]
Each $b_s(0)=1+h\mathbf k(s)\cdot\omega_0$ is positive. For $0<t\le1$, all imaginary parts have the same strict sign, so every $b_s(t)$ remains nonzero. The square roots continued from the positive real roots lie in the same right-hand quadrant. The identity principle identifies this expression with the holomorphic continuation of \eqref{curvedhomogeneous} for $0\le t<1$: the continued square root of $\xi\cdot\xi$ and its numerator roots are nonzero before the endpoint, and the equality holds initially on a real interval. The numerator roots have nonzero limits at $t=1$, and their average with the nonnegative weight $w$ also has a nonzero limit. The denominator is $z(t)\to0$, so the expression diverges, contradicting the bounded angular representation.

If the curvature vanishes identically on the interval, the curve is straight and \eqref{curvedhomogeneous} continues as a positive constant times $(1-t^2)^{-1/4}$, which also diverges. Thus positivity rules out rank-one focal patches. Together with the rank-zero case and the dichotomy $J\in\{0,1\}$, this proves Theorem~\ref{thm:nofocal}.
\end{proof}
\section{First-order normals along an arc}\label{app:arc}\label{sec:arc}
\begin{theorem}[First-order normals along an analytic arc]\label{thm:arcnormal}
Let $M\subset\R^3$ be an embedded topological surface with locally finite $\HH^2$, and let $\lambda$ be a positive measure with
$c_V\HH^2\restr M\le\lambda\le C_V\HH^2\restr M$ on each compact coordinate neighborhood $V$, where $0<c_V\le C_V<\infty$.
Let $\Gamma\subset M$ be a regular analytic open arc. Suppose an analytic ambient function $T$ satisfies $T=0$ on $\Gamma$, $dT\ne0$, and
\[
 (q,\pm dT(q))\notin\WF(\lambda)\qquad(q\in\Gamma).
\]
Then at $\HH^1$-almost every point of every smaller arc, $M$ has a nonzero Fr\'echet normal. More precisely, after analytic straightening to coordinates $(s,t=T,z)$ with $\Gamma=\{t=z=0\}$, almost every $b$ admits $\alpha_b\in\R$ such that
\begin{equation}\label{arcfirstorder}
 z-\alpha_b t=o\bigl(|(s-b,t,z)|\bigr)\quad\text{on }M\text{ as }(s,t,z)\to(b,0,0).
\end{equation}
\end{theorem}
\begin{proof}\leavevmode\par
\subsection{Transverse regularity and level geometry}

We prove Theorem~\ref{thm:arcnormal} by straightening the marked arc analytically. This change of coordinates is locally bi-Lipschitz, transports analytic wavefront, and preserves comparability of the measure with $\HH^2$. After shrinking the arc and its ambient neighborhood, closedness of wavefront gives a uniform missing double cone about $dt$.

\subsubsection{The geometric and analytic inputs}
We use the following geometric and analytic results.
\begin{enumerate}
\item Metric coarea: for a Lipschitz map $f:X\to\R$ on an arbitrary metric space and any $D\subset X$,
\begin{equation}\label{metriccoarea}
 \int_{\R}^{*}\HH^1(D\cap f^{-1}(r))\,dr
 \le \frac{4}{\pi}\operatorname{Lip}(f)\HH^2(D).
\end{equation}
Here $\int^{*}$ denotes the upper integral. If $X$ is boundedly compact, $D$ is $\HH^2$-measurable, and $\HH^2(D)<\infty$, the slice-measure function is measurable and the ordinary integral gives the same inequality \cite[Theorem 1.1]{EH}. All applications below take place on measurable subsets of a compact Euclidean coordinate neighborhood with finite area, so we use ordinary integration there.
\item A continuous function without local extrema on a topological surface is monotone in the local maximum--minimum sense. Every level component of a continuous monotone function on an open topological disk exits every compact subset of the disk \cite[Corollary 2.8]{Ntalam}. If the surface has locally finite area and the function is locally Lipschitz, almost every level is locally a proper rectifiable one-manifold \cite[Corollary 1.11]{EIR}; its metric escape property is given by \cite[Proposition 3.7]{EIR}.
\item A finite-length continuum is a locally connected arcwise connected continuum. A Peano continuum in the sphere that separates two specified points contains a Jordan curve separating those same points \cite[Lemmas 2.6--2.7]{EIR}. Local Jordan neighborhoods whose boundaries lie in a metric distance sphere are available \cite[Lemma 3.3]{EIR}.
\item Hausdorff limits of compact sets with at most $N$ connected components satisfy lower semicontinuity of $\HH^1$ \cite[Theorem 2.9]{AO}.
\item For a locally closed graph of a multifunction, the absence of a nonzero limiting normal of the form $(\tau,0)$ gives the Aubin property: nearby fibers move by a uniform Lipschitz bound in a fixed neighborhood \cite[Theorem 2.4 and Remark 2.1]{ML}.
\end{enumerate}
The $C^1$ supporting-function form of the wavefront theorem in Section~\ref{sec:frechet} places every Fr\'echet normal, and hence every limiting normal, in $\WF(\lambda)$. The missing $dt$ cone therefore implies that $t$ has no local extrema and gives, for
$A(t)=\{(s,z):(s,t,z)\in M\}$, fixed buffered neighborhoods and a constant $L$ with
\begin{equation}\label{fiberAubin}
 A(t_1)\cap W\subset A(t_2)+L|t_1-t_2|\overline B.
\end{equation}
Successive nearest-point choices on increasingly fine time partitions, followed by Arzel\`a--Ascoli and closedness, give an $L$-Lipschitz path in $(s,z)$ through each point, parametrized by $t$. We choose each path within the fixed buffered neighborhoods.

\subsubsection{Every nearby level has finite length}
\begin{lemma}\label{lem:alllevels}
Under the preceding assumptions every level $M\cap\{t=\tau\}$ has locally finite $\HH^1$.
\end{lemma}
\begin{proof}
For a nonnegative compact cutoff $\chi$, the ordinary wavefront gap gives a smooth density for the $t$-pushforward of $\chi\lambda$. Let $D$ bound this density on a smaller compact neighborhood. Applying \eqref{metriccoarea} to measurable time sets and using $\lambda\ge c_{\mathrm{src}}\HH^2$ gives
\[
 \HH^1(M\cap K_1\cap\{t=\tau\})\le \frac{4D}{\pi c_{\mathrm{src}}}
\]
for almost every $\tau$. Choose compact neighborhoods $K_0\subset\Int K\subset K\subset\Int K_1$ with positive separating distances. At such a generic time, trim the proper level curves at $\partial K_1$ and retain the compact connected pieces meeting $K$. The escape property ensures that every retained piece reaches $\partial K_1$ and thus has length at least $\delta=\operatorname{dist}(K,\partial K_1)>0$. Their number and total length are consequently uniformly bounded. For an exceptional time $\tau_0$, choose generic times $\tau_j\to\tau_0$. The Aubin inclusion approximates every point of $M\cap K_0\cap\{t=\tau_0\}$ by retained pieces. A subsequential Hausdorff limit contains this exceptional level in $K_0$, and the finite-component form of Go\l\k{a}b's theorem gives it finite length. Exhausting the neighborhood proves the assertion.
\end{proof}

\subsubsection{Finite forests and the marked zero arc}
\begin{lemma}\label{lem:forest}
Let $W$ be a Jordan disk compactly contained in a surface chart, let $t$ be continuous without local extrema, and suppose $Z=\overline W\cap\{t=0\}$ has finite length and meets $\partial W$ in a finite set $D$. Then $Z$ is a finite forest whose endpoints lie in $D$.
\end{lemma}
\begin{proof}
Every component $C$ of $Z$ meets $D$. If $p\in C\cap W$, the component of $\{t=0\}\cap W$ through $p$ lies in $C$ and exits every compact subset of $W$. Its connected closure in $\overline W$ remains in $C$ and meets $\partial W$, hence meets $D$. A component contained in the boundary already lies in $D$. Thus $Z$ has finitely many components, each a finite-length Peano continuum. A Jordan curve in $Z$ would enclose a disk on which a nonconstant $t$ attains a positive maximum or a negative minimum; if $t$ were constant on that disk, every interior point would be a local extremum. Each component is therefore a dendrite, with a unique arc between any two points. Let $F$ be the union of the finite trees joining the points of $D$ in each component.

The set $Z\cup\partial W$ is a Peano continuum: its finitely many finite-length components attach to the Jordan boundary, and a finite collection of surjective excursions gives a continuous parametrization. Every Jordan curve in this union lies in $F\cup\partial W$. Indeed, an interior arc between boundary points is the unique dendrite arc, while a loop meeting the boundary at most once would be a Jordan curve in a dendrite.

Let $V$ be a component of $W\setminus F$. If $V\setminus Z$ were disconnected, $Z\cup\partial W$ would separate two of its points. The Peano separator theorem would then give a Jordan curve separating those points. Such a curve lies in $F\cup\partial W$ and is disjoint from the connected set $V$, a contradiction. Hence $V\setminus Z$ is connected, and $t$ has one nonzero sign there. Any point of $Z\cap V$ would be a local extremum, so $Z\cap V$ is empty and $Z=F$. Near an interior endpoint of a finite forest, the punctured complement is connected and yields the same contradiction. All endpoints therefore lie on $\partial W$.
\end{proof}
One-dimensional coarea on a locally finite-length level, together with the Jordan-neighborhood result, gives each point a Jordan neighborhood whose boundary meets the level finitely many times. Lemma~\ref{lem:forest} applies in this neighborhood. A finite forest containing the marked embedded arc agrees locally with the arc away from its finitely many vertices. A countable chart cover therefore shows that the local zero fiber equals the marked arc at all but countably many points of $\Gamma$. Work in one such neighborhood. Comparison with time zero in \eqref{fiberAubin} gives
\begin{equation}\label{heightbound}
 |z|\le L|t|\quad\text{at every source point in a fixed smaller neighborhood.}
\end{equation}

\subsubsection{Jordan boxes at every scale}
\begin{lemma}\label{lem:levelboxes}
In an isolated-zero-arc neighborhood, two points of the same time in a fixed rescaled compact region can be joined by a compact level arc in a larger fixed rescaled compact region. The enlargement is independent of the sufficiently small scale. Moreover the projection $\pi(s,t,z)=(s,t)$ of the source covers a base rectangle.
\end{lemma}
\begin{proof}
Choose nested compact source neighborhoods with positive spatial buffers inside the fixed Aubin neighborhood. We first construct a coarse Jordan box with closure in the inner neighborhood and a time interval small enough to keep all backward paths in the larger neighborhood. Every fixed enlarged rescaled region will then fit within these buffers at sufficiently small scales. Choose zero-axis points on opposite sides of the origin and disjoint short Aubin paths through them. The zero arc separates a small intrinsic disk into two sides of opposite sign. Join the positive path endpoints by an embedded arc in the positive side avoiding the path interiors. Planar arc straightening shows that the two disjoint slit arcs from the disk boundary to free interior tips leave that side connected, so this connector exists. Construct a connector in the negative side in the same way. The four arcs form a Jordan boundary whose top and bottom are compact and have $t$ bounded away from zero. At every sufficiently small time $\tau$, the boundary meets the level in exactly two points. Lemmas~\ref{lem:alllevels} and~\ref{lem:forest}, together with a path joining an interior positive point to an interior negative point, show that the closed level inside the box is the unique arc joining these boundary points.

Fix $\kappa>0$ and $A>L\kappa$. At scale $r$, take Aubin paths through the zero-axis points $s=\pm Ar$ for $|t|\le\kappa r$. Their $s$ coordinates lie within $L\kappa r$ of $\pm Ar$, so the paths are disjoint. At each cap time $\pm\kappa r$, join their endpoints by the intervening subarc of the coarse box's level arc. Each cap meets the side paths only at its endpoints, and the two caps have different times. These four arcs form a fine Jordan boundary inside the coarse box. Its intersection with the coarse zero arc consists exactly of the two selected endpoints. The two tails of the coarse zero arc reach the coarse boundary without meeting the fine boundary and hence lie in its exterior. At each selected endpoint the side path crosses the zero arc, since its time parameter changes sign. Jordan separation therefore puts the open central interval $(-Ar,Ar)$ inside the fine box and the two tails outside. The zero level of the closed fine box is precisely $[-Ar,Ar]$. For each intermediate time, interior points sufficiently near the two caps have times on opposite sides of it, so an interior path between them meets the level. Lemmas~\ref{lem:alllevels} and~\ref{lem:forest}, with the two boundary intersections, give exactly one level arc.

Throughout the fine interior, $|t|<\kappa r$: a value beyond a cap time would force an interior maximum or minimum on the compact closure. Moreover,
\begin{equation}\label{boxbounds}
 |s|\le(A+2L\kappa)r,\qquad |z|\le L\kappa r.
\end{equation}
Indeed, suppose an interior point has $s>(A+2L\kappa)r$. Follow an Aubin path from its time to zero. Its spatial displacement is at most $L\kappa r$, so its $s$ coordinate stays above $(A+L\kappa)r$, whereas both side paths have $s\le(A+L\kappa)r$. The path's times lie strictly between the caps, and the fixed buffers keep it in the coordinate neighborhood. It therefore avoids the whole fine boundary. Jordan separation keeps this path in the fine interior, placing its endpoint in $[-Ar,Ar]$, which contradicts the same coordinate bound. The negative side is identical. The $z$ bound follows from \eqref{heightbound}.

Conversely, fix $M,T_0$, choose $\kappa>T_0$ and $A>M+L(T_0+\kappa)$, and consider a source point with $|s|\le Mr$, $|t|\le T_0r$. Along its Aubin path to time zero,
\[
 |s|\le (M+LT_0)r<(A-L\kappa)r.
\]
Thus the path misses both side paths and, since $|t|\le T_0r<\kappa r$, both caps. Its endpoint lies in the open central zero interval, already shown to be inside the fine box. Jordan separation then places the entire path, including its starting point, inside the box. Two such points at the same time lie on its unique level arc, and \eqref{boxbounds} gives a uniform enlargement. Finally, the endpoints of every sufficiently small fixed level have opposite $s$ coordinates. Connectedness supplies all intermediate coordinates, so the projection covers a base rectangle. Since $\pi$ is $1$-Lipschitz and $\lambda\ge c\HH^2$, we also obtain
\begin{equation}\label{basearealower}
 \pi_*\lambda\ge c\,ds\,dt
\end{equation}
on a smaller rectangle, with source restricted to the fixed neighborhood.
\end{proof}

\subsection{Positive moments, traces, and a common tangent measure}
\subsubsection{Division transverse to the arc}
\begin{lemma}\label{lem:division}
Suppose $V$ is a locally finite signed measure, the $t$-marginal of $|V|$ is locally absolutely continuous, and $U=t^kV$ has an ordinary wavefront gap in a cone about $dt$. Then $V$ has a gap in every fixed smaller such cone. The smaller cone may be chosen independently of the positive integer $k$.
\end{lemma}
\begin{proof}
Localize both measures with the same compact cutoff. Absolute continuity of the marginals and the Riemann--Lebesgue lemma imply that, for fixed spatial frequency $\sigma$, the Fourier transforms of $t^jV$, $0\le j\le k$, tend to zero as $|\tau|\to\infty$. Since
\[
 \partial_\tau^k\widehat V(\sigma,\tau)=(-i)^k\widehat U(\sigma,\tau),
\]
integration from $+\infty$ gives, for $\tau>0$,
\[
 \widehat V(\sigma,\tau)
 =\frac{i^k}{(k-1)!}\int_\tau^\infty
 (v-\tau)^{k-1}\widehat U(\sigma,v)\,dv.
\]
The preceding limits make all integration constants zero. On a smaller vertical cone, the entire integration ray stays in the original cone, so arbitrarily high rapid-decay bounds for $\widehat U$ give the corresponding bounds for $\widehat V$. For negative $\tau$, integrate from $-\infty$. Cutoff localization and the smaller-cone convolution estimate give the local assertion.
\end{proof}

By \eqref{heightbound}, we can choose a source cylinder whose vertical boundary bands are disjoint from $M$. Take a smooth compact cutoff $0\le\chi\le1$ equal to one on an ambient neighborhood of the selected smaller arc, using a compact base cutoff and a vertical cutoff whose transition bands miss $M$. Extending $\lambda^{\rm loc}=\chi\lambda$ by zero outside the cylinder gives a fixed finite positive compactly supported measure for all subsequent pushforwards. Put $a=z/t$ off $t=0$ and
\[
 V_k=\pi_*(\chi a^k\lambda),\qquad \omega=V_0.
\]
The zero arc has $\lambda$ mass zero, and $|V_k|\le L^k\omega$. The source gap gives smooth local $t$-marginals for $\omega$ and a common ordinary gap for
$t^kV_k=\pi_*(\chi z^k\lambda)$. Applying Lemma~\ref{lem:division} with these smooth localizations gives a common smaller ordinary gap for every $V_k$. We fix this cone once for all moment orders. The constants in its rapid-decay estimates may depend on $k$; only the cone aperture is common.

\subsubsection{Positive traces and area growth}
\begin{lemma}\label{lem:positivetrace}
A positive compact measure $V$ on $(s,t)$ with a uniform ordinary gap about $dt$ has a positive Radon trace $m$ on $t=0$. At a point $b$ where $m$ has finite Lebesgue density $\ell$ and its singular part has density zero,
\[
 V(B_r(b,0))=O(r^2).
\]
If $\ell=0$, the conclusion is $o(r^2)$. The same statements hold after localizing a locally finite measure.
\end{lemma}
\begin{proof}
The wavefront gap permits noncharacteristic distributional restriction. For a nonnegative test $\varphi(s)$, the $t$-pushforward of $\varphi V$ has a smooth nonnegative density whose value at zero is the trace on $\varphi$. A larger cutoff bounds this functional on compact spatial supports, so the trace is a positive distribution of order zero.

For the growth estimate, set $P_y(x)=y/[\pi(x^2+y^2)]$ and choose $c$ so large that the complement of the missing cone lies in $|\tau|<c|\sigma|$. An angular Fourier cutoff, with low frequencies assigned to the remainder, gives
\begin{equation}\label{angularsplit}
 V=v+w\,ds\,dt,
\end{equation}
where $w$ is Schwartz and $v=V-w\,ds\,dt$ is a finite signed measure whose Fourier support lies in a strict horizontal double cone away from the origin. Derivatives of $\widehat V$ are transforms of the compact measure multiplied by coordinate polynomials, so they satisfy the same rapid cone estimates. The complementary Fourier part, including low frequencies, therefore has a Schwartz inverse. In coordinates $x_\pm=s\pm ct$, the Fourier variables have the same sign on this cone, and their product-Poisson multiplier is $e^{-y|\sigma|}$. Properness of $(\sigma,\tau)\mapsto\sigma$ on the cone and boundedness of $\widehat v$ justify Fourier integration and yield
\begin{equation}\label{poissonsplit}
 \int K_{y,b}\,dV
 =P_y*m(b)+\int K_{y,b}w\,ds\,dt-P_y*w(\cdot,0)(b),
 \quad K_{y,b}=2cP_y(s-b+ct)P_y(s-b-ct).
\end{equation}
The last two terms cancel as $y\downarrow0$, while $P_y*m(b)\to\ell$ at the stated density point. On $B_y(b,0)$, $K_{y,b}\ge c_1y^{-2}$. Positivity now gives the $O(y^2)$ bound and, when $\ell=0$, the $o(y^2)$ bound.
\end{proof}

For a polynomial $p\ge0$ on $[-L,L]$, the trace $m_p$ of $\pi_*(\chi p(a)\lambda)$ is positive, and
$|m_p|\le\|p\|_{[-L,L]}m_0$ for arbitrary real $p$. Equation~\eqref{basearealower} gives $m_0\ge c\,ds$. Choose a simultaneous Lebesgue differentiation point for the absolutely continuous densities of the countably many moment traces and their variations, at which the singular variations have density zero. Such points form a full-measure set, and the rescaled traces there converge to constant Lebesgue densities. Denote these densities by $\ell_k$, with $\ell_0\ge c>0$. Positivity on rational polynomials and uniform approximation determine a unique finite positive law $\nu$ on $[-L,L]$ satisfying
\begin{equation}\label{lawdefinition}
 \int a^k\,d\nu(a)=\ell_k.
\end{equation}

\subsubsection{Spectral support and traces of the tangent measure}
Translate the selected point to zero and dilate by $D_r x=x/r$. Lemma~\ref{lem:positivetrace} and bounded slopes give locally bounded rescaled ambient and base measures
\[
 r^{-2}(D_r)_*(\chi\lambda),\qquad r^{-2}(D_r)_*V_k.
\]
Choose a common subsequence $r_j\downarrow0$ along which these measures and the joint positive measures obtained from $\chi\lambda$ under $(s,t,z)\mapsto(s/r_j,t/r_j,z/t)$, with mass factor $r_j^{-2}$, converge locally. Write $W_{k,j}=r_j^{-2}(D_{r_j})_*V_k$, denote the ambient limit by $\lambda_\infty$, and denote the base moment limits by $W_k$. Compactness of $[-L,L]$ identifies all moments simultaneously. The local quadratic bound controls balls for which $r_jR$ is small, and the finite total mass of $\chi\lambda$ controls the remaining radii. Thus the rescaled measures are uniformly tempered, and their limits have global quadratic growth about zero.

\begin{lemma}\label{lem:tangentcone}
For one constant $C_0$ independent of $k$,
\begin{equation}\label{tangentcone}
 \supp\widehat W_k\subset\{ |\tau|\le C_0|\sigma|\},
 \qquad W_k|_{t=0}=\ell_k\,ds.
\end{equation}
The base limit has no mass on $t=0$, and
$W_k=\pi_*(a^k\lambda_\infty)$ with $a=z/t$.
\end{lemma}
\begin{proof}
For each fixed $k$, the growth bounds from Lemma~\ref{lem:positivetrace} and $|V_k|\le L^k\omega$ give uniform tempered-distribution bounds for $W_{k,j}$. Fix a closed double cone strictly inside the common missing cone supplied by Lemma~\ref{lem:division}. On it the compact source moments satisfy
\[
 |\widehat V_k(\xi)|\le C_{k,N}(1+|\xi|)^{-N}
 \quad\text{for every }N,
 \qquad
 \widehat W_{k,j}(\xi)=r_j^{-2}\widehat V_k(\xi/r_j).
\]
For a frequency test compactly supported in this cone away from zero, choosing $N>2$ makes the rescaled pairing tend to zero. The limiting Fourier transforms therefore vanish in that cone. Its fixed aperture supplies one constant $C_0$ for the support assertion, independent of $k$.

For trace convergence, fix a smooth cutoff $\rho(s,t)$ compactly supported in the scaled variables and equal to one near the compact axis interval under consideration. Compact local mass bounds give a uniform bound on $\widehat{\rho W_{k,j}}$ for each fixed $k$. To obtain a uniform bound on its vertical-frequency tails, choose a compact frequency cutoff $\theta$ equal to one on the unit ball and an angular cutoff $\beta$ supported strictly in the common missing cone and equal to one on a smaller vertical cone. Before multiplication by $\rho$, decompose the transform as
\[
 \underbrace{\theta\widehat W_{k,j}}_{\text{compact low frequency}},
 \qquad
 \underbrace{(1-\theta)\beta\widehat W_{k,j}}_{\text{high vertical frequency}},
 \qquad
 \underbrace{(1-\theta)(1-\beta)\widehat W_{k,j}}_{\text{angular complement}}.
\]
All the following constants may depend on $k$, but are uniform in $j$. On $|\xi|\ge1$, the high vertical part is bounded by
\[
 C_{k,N}r_j^{-2}(1+|\xi|/r_j)^{-N}
 \le C_{k,N}r_j^{N-2}|\xi|^{-N}.
\]
Choosing $N$ sufficiently large gives uniformly bounded weighted $L^1$ norms of every prescribed order. Convolution with $\widehat\rho$ preserves rapid decay on a smaller cone. The angular complement has uniformly bounded tempered seminorms, and its frequency support is angularly separated from a still smaller vertical cone. Pairing this part with translates of the Schwartz function $\widehat\rho$ therefore gives uniform rapid decay on that cone. Finally, $\theta\widehat W_{k,j}$ has compact support and uniformly bounded distribution seminorms, so its convolution with $\widehat\rho$ is uniformly rapidly decreasing at large frequency. Thus $\widehat{\rho W_{k,j}}$ has uniform rapid decay on one fixed vertical cone, whose aperture is independent of $k$.

In the Fourier formula for restriction to $t=0$, pair $\sigma$ with the Fourier transform of a compactly supported smooth spatial test. On the vertical cone the preceding decay is uniformly integrable in $\tau$. On its complement the transforms are uniformly bounded, and for each $\sigma$ the allowed $\tau$ interval has length $O(1+|\sigma|)$, which is absorbed by the spatial test's Schwartz decay. These two bounds give an integrable majorant independent of $j$. Weak convergence of the localized measures gives pointwise convergence of their transforms, so dominated convergence gives convergence of the traces. With $D_r^s(s)=s/r$, the original traces rescale as
\[
 W_{k,j}|_{t=0}=r_j^{-1}(D_{r_j}^s)_*m_k
 \longrightarrow \ell_k\,ds
\]
by differentiation at the chosen point. This proves the trace assertion.

The exact conic Fourier support also gives an analytic wavefront gap outside the cone. To see this, use the Gaussian-window FBI transform
\[
 \mathcal T_\varepsilon u(x,\xi)
 =\left\langle u(y),
 e^{i(x-y)\cdot\xi/\varepsilon-|x-y|^2/(2\varepsilon)}
 \right\rangle.
\]
Its Fourier representation pairs $\widehat u(\eta)$, up to a harmless normalization and phase, with
$e^{ix\cdot\eta}e^{-\varepsilon|\eta-\xi/\varepsilon|^2/2}$.
For $\xi$ in a compact cone disjoint from the spectral support and $x$ in a compact set, spectral separation gives an $e^{-c/\varepsilon}$ bound multiplied by a polynomial tempered-seminorm factor, which can be absorbed into the exponential. Inserting a physical cutoff equal to one near $x$ changes the pairing by a Gaussian tail with the same exponential bound. The local Gaussian FBI characterization \cite[Appendix B, Proposition B.1]{SW} therefore excludes the corresponding analytic wavefront covectors. In particular, every localized positive $W_0$ has a smooth $t$-marginal and hence zero axis mass. The joint slope limit is dominated by $W_0$, so it also has zero axis mass. Off the axis, the continuous relation $z=ta$ passes to the joint limit and identifies the moments as claimed.
\end{proof}

\subsection{Product-Poisson identities for the tangent measure}
\begin{lemma}\label{lem:tangentpoisson}
With $c>C_0$ and $K_{y,b}$ as in \eqref{poissonsplit}, the tangent measure satisfies
\begin{equation}\label{allpoisson}
 \int K_{y,b}(s,t)\,g(z/t)\,d\lambda_\infty
 =\int g\,d\nu\qquad(y>0,\ b\in\R)
\end{equation}
for every continuous $g$ on $[-L,L]$. Every such weighted base measure has the same exact Fourier cone.
\end{lemma}
\begin{proof}
Using Lemmas~\ref{lem:positivetrace} and~\ref{lem:tangentcone}, begin with a nonnegative polynomial $p$. Its compact source measure $V$ is positive and has trace $m$ with density $\ell=\ell_p$ at zero. For $V_j=r_j^{-2}(D_{r_j})_*V$, \eqref{poissonsplit} gives
\[
 \int K_{y,b}\,dV_j=P_y*m_j(b)
 +\int K_{y,b}w(r_js,r_jt)\,ds\,dt
 -P_y*w(r_j\cdot,0)(b),
\]
where $m_j=r_j^{-1}(D_{r_j})_*m$. The finite density at zero and finite total mass of $m$ give the uniform global bound $m_j([-R,R])\le AR$. The Poisson tails are therefore uniformly $O(R^{-1})$, and local trace convergence gives $P_y*m_j\to\ell$. The two $w$ terms cancel by bounded convergence, so the displayed integrals tend to $\ell$. Positivity and Fatou's lemma yield, for $\Omega=W_p$,
\begin{equation}\label{poissonfatou}
 0\le U(b,y):=\int K_{y,b}\,d\Omega\le\ell.
\end{equation}
In particular, $K_{1,0}\Omega$ is a finite measure.

Let $\widetilde\Omega$ be the pushforward of $\Omega$ under $x_\pm=s\pm ct$. The dual coordinates are $\eta_\pm=(\sigma\pm\tau/c)/2$. Since $c>C_0$, the exact Fourier cone is contained in $[0,\infty)^2\cup(-\infty,0]^2$, including the origin. Moreover,
\[
 \int\frac{d\widetilde\Omega(x_+,x_-)}{(1+x_+^2)(1+x_-^2)}
 =\frac{\pi^2}{2c}\int K_{1,0}\,d\Omega<\infty.
\]
The Fourier characterization \cite[Theorem 3.2]{NS} therefore makes $\widetilde\Omega$ a Nevanlinna measure. Its product-Poisson extension
\[
 F(w_+,w_-)=\int P_{\operatorname{Im}w_+}(x_+-\operatorname{Re}w_+)
 P_{\operatorname{Im}w_-}(x_--\operatorname{Re}w_-)\,d\widetilde\Omega
\]
is pluriharmonic on the product of upper half-planes \cite[Section 2.1]{NS}. Restriction to the holomorphic diagonal $w\mapsto(w,w)$ gives $U(b,y)=2cF(b+iy,b+iy)$, which is harmonic on the upper half-plane.

To identify its boundary value, choose $0\le\chi\le1$ compact and equal to one near the axis support of a spatial test $\varphi$. Formula \eqref{poissonsplit} applied to $\chi\Omega$ gives the local trace limit. The remaining tail satisfies
\begin{equation}\label{kerneltail}
 (1-\chi(s,t))\int |\varphi(b)|K_{y,b}(s,t)\,db
 \le C_{\varphi,\chi}y K_{1,0}(s,t),\quad 0<y<1.
\end{equation}
In coordinates $x_\pm$, at least one coordinate outside the cutoff neighborhood is separated by a fixed distance from $b\in\supp\varphi$. Its Poisson factor is bounded by $Cy/(1+x_\pm^2)$. If the other coordinate is bounded, integrate its Poisson factor in $b$ and use the positive lower bound for its product-Cauchy weight. If it is unbounded, it supplies a second factor $Cy/(1+x_\mp^2)$. This proves \eqref{kerneltail}. Integrability bounds the tail by $O(y)$, so \eqref{tangentcone} gives the distributional boundary value $\ell ds$. By \eqref{poissonfatou}, $U$ is a bounded harmonic function with this constant boundary value and is therefore identically $\ell$. This proves \eqref{allpoisson} for nonnegative polynomials, and hence for all polynomials.

The case $p=1$ gives finite product-kernel mass. Uniform polynomial approximation on the compact slope interval then proves \eqref{allpoisson} for every continuous $g$. Domination by the quadratic-growth base measure makes the same approximation valid in the space of tempered measures, extending the common exact Fourier cone to every $g$.
\end{proof}

\subsection{Geometry of the tangent measure}
\begin{lemma}\label{lem:limitgeometry}
The support $M_\infty=\supp\lambda_\infty$ has locally finite $\HH^2$ and, on each compact neighborhood $V$, satisfies $\lambda_\infty\ge c_{\infty,V}\HH^2\restr M_\infty$ for a positive constant $c_{\infty,V}$ associated with the limit. It also satisfies $|z|\le L|t|$ and the Aubin comparison. Every exact-time fiber is connected and projects onto the whole $s$-axis; two points in a bounded region of one fiber can be joined by a continuum in a larger bounded region.
\end{lemma}
\begin{proof}
We retain the subsequence of Lemma~\ref{lem:tangentcone} and write $\lambda_j=r_j^{-2}(D_{r_j})_*(\chi\lambda)$. A uniform lower area estimate holds at every source point in a fixed neighborhood where $\chi=1$. For such a point $p$ and a small radius $r$, Aubin motion supplies points within $r/4$ of $p$ at all times $|\tau-t(p)|<c_1r$. At almost every such time, the proper level component through the corresponding point exits a fixed compact ball in the surface chart. Its portion before leaving $B_r(p)$ has length at least $r/2$. Coarea gives $\lambda(B_r(p))\ge c_2r^2$, with constants uniform on the source neighborhood. Every fixed rescaled compact region eventually lies in this $\chi=1$ neighborhood, so the rescaled localized measures inherit the estimate.

Suppose source support points $p_j$ converge to $p$ in a fixed rescaled compact region. For each sufficiently small $r>0$, the lower estimate and $p_j\to p$ give
\[
 \lambda_j(B_{r/2}(p_j))\ge c_2r^2/4,
 \qquad B_{r/2}(p_j)\subset\overline B_{3r/4}(p)
\]
for all large $j$. The compact-set Portmanteau inequality yields
\[
 \lambda_\infty(B_r(p))
 \ge\lambda_\infty(\overline B_{3r/4}(p))
 \ge\limsup_j\lambda_j(\overline B_{3r/4}(p))
 \ge c_2r^2/4.
\]
Thus every Hausdorff limit point of the source supports lies in $\supp\lambda_\infty$ and has a positive lower ball bound. Conversely, weak convergence places the limit support in the local Hausdorff limit of the source supports. The disjoint-ball covering argument gives $\HH^2\restr(M_\infty\cap V)\le c_{\infty,V}^{-1}\lambda_\infty\restr V$ on compact neighborhoods, proving local finiteness of area. Closedness passes the height bound and Aubin comparison to the limit. For exact-time connectivity, first use Aubin motion to adjust approximating source points to a common time. Lemma~\ref{lem:levelboxes} then joins any two of these points in a fixed enlarged rescaled compact set, and a Hausdorff limit of the connecting arcs is a continuum in the exact-time limit fiber. Applying the same boxes with arbitrarily separated endpoints shows that its spatial projection is all of $\R$. The geometric Hausdorff limits are taken along refinements of the moment subsequence, so all established identities continue to hold.
\end{proof}

\subsection{Classification of the compact slope law}
\begin{proposition}\label{prop:lawdirac}
The law $\nu$ in \eqref{lawdefinition} is $\ell_0\delta_\alpha$ for some $\alpha\in[-L,L]$.
\end{proposition}
\begin{proof}\leavevmode\par
\subsubsection{Slope variation yields an absolutely continuous component}
Suppose a fiber at $t=T\ne0$ contains two distinct heights. Aubin motion preserves the separated slope values on a nearby time interval. By Lemma~\ref{lem:limitgeometry}, connecting continua remain in a common compact region and contain every intermediate slope. For each $r$ in an open slope interval, the set $\{a=r\}$ in that region therefore projects onto a fixed time interval and has a uniform positive lower bound for $\HH^1$. The function $a=z/t$ is Lipschitz on this compact region. Applying \eqref{metriccoarea} to $a^{-1}(D)$ for measurable slope sets $D$, together with the lower area weight and a positive lower bound for $K_{1,0}$, gives
\[
 \nu(D)\ge c\,|D|
\]
on the slope interval. Therefore $\nu$ has a nonzero absolutely continuous part.

If every fiber has a single height, its full spatial projection gives $M_\infty=\{z=g(t)\}$, and the Aubin comparison makes $g$ locally Lipschitz. A nonconstant $g(t)/t$ on either open half-line gives the same conclusion: choose a compact time corridor whose slopes cover an interval, and apply coarea to its product with a compact spatial interval. It follows that a slope law with zero absolutely continuous part has a constant slope on each time half-space.

\subsubsection{Normalized conditionals and their exact trace}
Assume now that $d\nu_{\rm ac}/dr=v(r)>0$ on a set of positive measure. Disintegrate the finite measure $K_{1,0}\lambda_\infty$ under $a=z/t$ as $\gamma_r\,d\nu(r)$, with each $\gamma_r$ a probability measure. Set
\[
 \lambda_r=K_{1,0}^{-1}\gamma_r,\qquad \omega_r=\pi_*\lambda_r.
\]
These measures are tempered, since $K_{1,0}^{-1}$ is bounded by a polynomial. Choose a countable dense family of smooth Fourier tests compactly supported outside the cone in \eqref{tangentcone} (Lemma~\ref{lem:tangentcone}), and a countable dense set of parameters $(y,b)\in(0,\infty)\times\R$. Lemma~\ref{lem:tangentpoisson}, tested against arbitrary continuous slope functions, gives the Fourier vanishing and \eqref{allpoisson} for the conditionals outside one $\nu$-null set for each selected test. Removing their countable union gives a common full-measure set. Continuity of each tempered distribution extends the Fourier identities to all tests outside the cone. On compact parameter sets, $K_{y,b}/K_{1,0}$ is uniformly bounded, so dominated convergence extends the kernel identities to all $y>0$ and $b\in\R$. On this common set,
\begin{equation}\label{conditionalidentities}
 \supp\widehat\omega_r\subset\{|\tau|\le C_0|\sigma|\},
 \quad \int K_{y,b}\,d\omega_r=1,
 \quad \omega_r|_{t=0}=ds.
\end{equation}
The trace equality follows from \eqref{poissonsplit} on compact localizations and the tail estimate \eqref{kerneltail}, using $\int K_{1,0}\,d\omega_r=1$. Indeed, the kernel identities give the constant distributional boundary value $ds$, and the tail estimate identifies it with the local trace. This argument applies first to a countable dense family of spatial tests and then to every test by continuity. Positivity and the ordinary gap also give zero axis mass.

We next establish local finite length of the conditional support. On a countable compact exhaustion $V_j$ of $\{t\ne0\}$, the function $a=z/t$ is Lipschitz and $\HH^2(M_\infty\cap V_j)<\infty$. The metric coarea inequality \eqref{metriccoarea} gives
\[
 \HH^1\bigl(M_\infty\cap V_j\cap\{a=r\}\bigr)<\infty
 \quad\hbox{for Lebesgue-almost every }r,
\]
simultaneously for all $j$. The defining slope map and a countable open basis also show that, for $\nu$-almost every $r$, $\lambda_r$ is carried by the closed set $M_\infty\cap\{z=rt\}$ and has zero axis mass. Every $\nu$-null exceptional set is Lebesgue-null on $\{0<v(r)<\infty\}$. We can therefore choose an $r$ in this set satisfying all the identities and finite-length conclusions.

On the plane $\{z=rt\}$, the projection $\pi(s,t,z)=(s,t)$ is a proper bi-Lipschitz bijection. Thus, for $E_r=\supp\omega_r$,
\begin{equation}\label{conditionalsupportlength}
 \HH^1(E_r\cap U)<\infty
 \quad\hbox{on every compact neighborhood }U\subset\{t\ne0\}.
\end{equation}
Indeed, its inverse image in that plane is compact and lies in one member of the preceding exhaustion, so the finite-length estimate applies directly to the support.

\subsubsection{Analytic branches from finite support length}
Fix $p=(s_0,t_0)\in E_r$ with $t_0\ne0$. The exact cone in \eqref{conditionalidentities} gives a uniform analytic gap about both signs of $dt$ near $p$. Finite length in a compact neighborhood implies that only countably many lines through $p$ meet $E_r$ there in positive $\HH^1$ measure: their intersections are disjoint except for $p$, and a finite measure space has at most countably many disjoint positive-measure subsets. Choose a sufficiently small $\varepsilon$ outside this exceptional set and introduce coordinates
\[
 S=s,\qquad T=t+\varepsilon(s-s_0),
\]
so that $\pm dT$ stay in the missing cones and the initial fiber $T=t_0$ is length-null.

Choose spatial endpoints outside this closed null fiber on either side of $p$. Closedness gives support-free bands about them over a smaller $T$ interval $I$. Let $\widetilde\omega_r$ be the measure in the new coordinates, and choose $\theta(S)$ equal to one between the bands and zero beyond their outer edges, with all transitions inside the bands. The cutoff is locally constant near the support and thus preserves the analytic wavefront gap. Projection onto $T$ is proper on the selected support, so analytic pushforward gives
\begin{equation}\label{analyticatomicmoments}
 q_k(T)\,dT=(\pi_T)_*(S^k\theta\widetilde\omega_r),\qquad k\ge0,
\end{equation}
with all $q_k$ real analytic on $I$.

The positive marginal $q_0\,dT$ gives a disintegration of the selected measure as positive compactly supported spatial measures $\eta_T\,dT$. On a smaller interval with compact closure, \eqref{conditionalsupportlength} and one-dimensional metric coarea \cite[Theorem 1.1]{EH} imply that the selected support fibers are finite for almost every $T$. The $\eta_T$ are therefore finite atomic measures almost everywhere. Choose a common full-measure set for this property and the countably many moment identities.

Put $H_m(T)=(q_{i+j}(T))_{0\le i,j<m}$. At each time in the common full-measure set, this is the Gram matrix of $1,S,\ldots,S^{m-1}$ in $L^2(\eta_T)$, and its determinant is positive precisely when $\eta_T$ has at least $m$ distinct atoms. If $\det H_m$ were not identically zero for arbitrarily large $m$, then outside the countable union of their discrete zero sets all these matrices would be invertible, contradicting finite atomicity on the common full-measure set. Since $p$ lies in the support, the selected measure is nonzero; thus a largest such $m\ge1$ exists. Positivity shows that the conditional measure has exactly $m$ distinct atoms at almost every time outside the zero set of $\det H_m$.

The Hankel recurrence gives the monic degree-$m$ polynomial whose roots are the atom positions. Its meromorphic coefficients are bounded almost everywhere because the roots stay in a compact spatial interval. Their poles are therefore removable, and continuity makes the extended polynomial real-rooted at every real time. Rellich's one-parameter theorem \cite[Theorem 3.4]{KP} gives complete analytic roots $S=b_1(T),\ldots,b_m(T)$, generically distinct. The Vandermonde formulas give meromorphic weights $W_j$ that are nonnegative and bounded almost everywhere by $q_0$. Their poles are also removable, and
\[
 \theta\widetilde\omega_r=
 \sum_{j=1}^m W_j(T)\,\delta_{S=b_j(T)}\,dT
\]
holds as an equality of measures, since absolute continuity of the time marginal gives zero mass at exceptional times.

Retain the branches with nonzero analytic weights. Each weight is positive at arbitrarily close noncollision times, so closedness places its complete graph in the support. Conversely, the finite union of these closed graphs carries the measure and is therefore exactly the selected support. The weights are positive except possibly at isolated times.

At a noncollision point with nonzero weight, the analytic curve's conormal belongs to analytic wavefront. Closedness extends this conclusion to weight zeros and collisions. In the original coordinates, a branch has the form
\[
 s=b_j(T),\qquad t=T-\varepsilon(b_j(T)-s_0).
\]
A zero of $1-\varepsilon b_j'(T)$ would make $dt$ a conormal, contradicting the wavefront gap. Original time is therefore a noncritical analytic parameter on every retained branch. After shrinking around $p$ to retain the branches through it, analytic inversion gives a finite family of complete graphs $s=\beta_j(t)$ with nonnegative analytic weights, possibly vanishing at isolated times. Applied at each support point, this construction proves local finiteness of the fibers and continuation of every support branch through interior times.

\subsubsection{Bounded branch speed and the axis trace}
At a nonzero-weight point of an isolated branch $s=\beta(t)$, the conormal $(1,-\beta'(t))$ belongs to analytic wavefront. Equation~\eqref{conditionalidentities} gives
\begin{equation}\label{branchspeed}
 |\beta'(t)|\le C_0.
\end{equation}
Closedness and continuity extend this bound through isolated weight zeros and collisions. The following continuation lemma applies to the support just constructed.

\begin{lemma}[Finite continuation in a bounded strip]\label{lem:finitecontinuation}
Let $E$ be relatively closed in $\R\times\{t\ne0\}$. Suppose that near each of its points, $E$ is a finite union of complete real analytic graphs $s=\beta_j(t)$, with identical germs identified, and every graph satisfies $|\beta_j'|\le C_0$. One can choose $T>0$ such that, for every $R>0$, the portion of $E$ with $|s|\le R$ and $0<|t|<T$ is covered by finitely many analytic branches continuing to $t=T$ or $t=-T$, respectively. Each branch continues toward the axis and has a unique limit there.
\end{lemma}
\begin{proof}
Take a countable cover by the local finite-branch neighborhoods. The collision times of nonidentical branches form a countable set, since each difference is analytic and has isolated zeros. Choose $T>0$ such that neither $T$ nor $-T$ is a collision time. At those times every support point has a unique analytic support germ.

Fix $R>0$ and a branch point with $|s|\le R$ and $0<t<T$. Its continuation toward $T$ stays in $|s|\le R+C_0T$. If it ended at an interior time $\tau$, bounded speed would give a spatial limit at $\tau$, belonging to $E$ by relative closedness. Write the finitely many local graphs there as $s=\beta_1(t),\ldots,\beta_m(t)$. The incoming graph $s=\beta(t)$ satisfies $\prod_{j=1}^m(\beta(t)-\beta_j(t))=0$ on a one-sided interval. The analytic identity theorem makes one factor identically zero there, so that local graph continues the incoming germ through $\tau$. Uniqueness of analytic continuation also fixes the germ through a crossing. Thus the branch reaches $T$.

Local finiteness of the fibers and relative closedness imply that $E\cap\{t=T,\ |s|\le R+C_0T\}$ is finite. Every initial branch continues to one of its finitely many germs, and uniqueness identifies it with the continuation of that germ. Continuing each terminal germ backward cannot end at a positive interior time, by the same argument. The speed bound gives a unique finite spatial limit as $t\downarrow0$. Apply the argument on the negative half-plane to obtain the remaining finitely many branches and limits.
\end{proof}

Apply Lemma~\ref{lem:finitecontinuation} to $E_r$. For a fixed $R>0$, let the finite set of axis limits be those of the branches covering $|s|\le R$, $0<|t|<T$. Choose a nonempty open interval $I$ whose closure lies in $(-R,R)$ and avoids these limits. Finiteness and convergence of the branches give $\varepsilon>0$ such that $I\times(-\varepsilon,\varepsilon)$ has no off-axis support. Since the axis has zero mass, its trace vanishes on $I$, contradicting the Lebesgue trace in \eqref{conditionalidentities}. Thus $\nu$ has zero absolutely continuous part.

\subsubsection{Completion of Proposition~\ref{prop:lawdirac}}
The slope-variation argument now gives a single height $g(t)$ per fiber and a constant slope on each open half-space, say $b_+$ and $b_-$. If these values differ, choose a nonnegative continuous function of the slope that selects one value. The resulting positive weighted base measure is supported on exactly one closed time half-plane, with full support up to its boundary. Its vertical support normal belongs to analytic wavefront, contradicting the common exact cone. Hence $b_+=b_-=: \alpha$, and \eqref{allpoisson} gives $\nu=\ell_0\delta_\alpha$.
\end{proof}
\subsection{From the tangent law to a first-order normal}
\begin{lemma}\label{lem:variance}
At the selected original source point the identity $\nu=\ell_0\delta_\alpha$ implies \eqref{arcfirstorder}.
\end{lemma}
\begin{proof}
Using the law from Proposition~\ref{prop:lawdirac}, the positive projected variance measure of the original source
$V=\pi_*(\chi(z/t-\alpha)^2\lambda)$ has trace density zero. Lemma~\ref{lem:positivetrace} gives $V(B_{Rr})=o(r^2)$ for each fixed $R$. For sufficiently small $r$, the cutoff is one on the ambient ball $B_{Rr}$, and its projection lies in the corresponding base ball. Positivity therefore gives
\begin{equation}\label{variancevanishes}
 \int_{M\cap B_{Rr}}(z/t-\alpha)^2\,d\lambda=o(r^2).
\end{equation}
Suppose points $p_j$ satisfy $|p_j|\le A|t_j|$ and $|z_j-\alpha t_j|\ge\varepsilon|t_j|$, with $t_j\to0$. Put $r=|t_j|$ and pass to one sign of $t_j$. Aubin motion preserves a point with height difference at least $7\varepsilon r/8$ over a time interval $J_j$ of length $c\varepsilon r$, while keeping these points in a fixed enlarged ball. Choose $c$ small enough that $r/2\le|t|\le3r/2$ throughout $J_j$. Choose the enlarged ball $B_{Rr}$ with radial margin at least $r$ around the moving points and contained in the chart for small $r$. At almost every time in $J_j$, the proper level component through the corresponding point exits the ball. Follow it until its height difference first reaches $\varepsilon r/2$ or it first exits $B_{Rr}$. Along this portion the difference is at least $\varepsilon r/2$. Reaching the lower difference changes $z$ by at least $3\varepsilon r/8$; exiting the ball traverses the radial margin. The portion therefore has length at least $\min(3\varepsilon/8,1)r$.

Apply \eqref{metriccoarea} to the Borel set
\[
 D_j=M\cap B_{Rr}\cap\{t\in J_j\}
       \cap\{|z-\alpha t|\ge\varepsilon r/2\}.
\]
The preceding portions give $\HH^1(D_j\cap\{t=\tau\})\ge c_1r$ for almost every $\tau\in J_j$. Coarea and the lower area weight yield $\lambda(D_j)\ge c_2r^2$. On $D_j$, the variance factor is at least $\varepsilon^2/9$, because $|t|\le3r/2$. Hence the variance mass in $B_{Rr}$ is bounded below by a positive multiple of $r^2$, contradicting \eqref{variancevanishes}.

Thus $(z-\alpha t)/|t|\to0$ in every cone $|p|\le A|t|$. Comparison with the zero axis also gives $|z-\alpha t|\le(L+|\alpha|)|t|$. For approaches with $|t|\ge\eta|p|$, use the conical limit; for those with $|t|<\eta|p|$, use this bound and let $\eta\downarrow0$. This proves \eqref{arcfirstorder}.
\end{proof}
The simultaneous density points form a full-measure set in each isolated-zero-arc chart. Removing the countably many forest vertices and applying a countable cover preserves this full-measure conclusion. Lemma~\ref{lem:variance} therefore completes the proof of Theorem~\ref{thm:arcnormal}.
\end{proof}
\begin{thebibliography}{99}
\bibitem{AO} G. Alberti and M. Ottolini, \emph{On the structure of continua with finite length and Go\l\k{a}b's semicontinuity theorem}, arXiv:1705.09941. \url{https://arxiv.org/pdf/1705.09941}.
\bibitem{BDH} C. Brouder, N. V. Dang, and F. H\'elein, \emph{Continuity of the fundamental operations on distributions having a specified wave front set (with a counter example by Semyon Alesker)}, arXiv:1409.7662. \url{https://arxiv.org/pdf/1409.7662}.
\bibitem{DP} A. Daniilidis and C. H. J. Pang, \emph{Continuity and differentiability of set-valued maps revisited in the light of tame geometry}, author manuscript. \url{https://www.arisdaniilidis.at/DP_2010.pdf}.
\bibitem{EH} B. Esmayli and P. Haj\l asz, \emph{The coarea inequality}, arXiv:2006.00419. \url{https://arxiv.org/pdf/2006.00419}.
\bibitem{EIR} B. Esmayli, T. Ikonen, and K. Rajala, \emph{Coarea inequality for monotone functions on metric surfaces}, author manuscript. \url{https://users.jyu.fi/~kirajala/publications/EIR_Coarea_Inequality_Monotone_v3.pdf}.
\bibitem{Ghomi} M. Ghomi, \emph{Open problems in geometry of curves and surfaces}. \url{https://ghomi.math.gatech.edu/Papers/op.pdf}.
\bibitem{GHL} M. Ghomi, R. Howard, and M. Lai, \emph{A converse to the Archimedes sphere theorem}, preprint, revised October 3, 2026. \url{https://ghomi.math.gatech.edu/Papers/archimedes.pdf}.
\bibitem{GT} D. Gilbarg and N. S. Trudinger, \emph{Elliptic Partial Differential Equations of Second Order}, Classics in Mathematics, Springer, 2001. \url{https://doi.org/10.1007/978-3-642-61798-0}.
\bibitem{HormSP} L. H\"ormander, \emph{The Analysis of Linear Partial Differential Operators I: Distribution Theory and Fourier Analysis}, second edition, Springer, 2003. \url{https://doi.org/10.1007/978-3-642-61497-2}.
\bibitem{Horm} L. H\"ormander, \emph{Remarks on Holmgren's uniqueness theorem}, Ann. Inst. Fourier \textbf{43}(5) (1993), 1223--1251. \url{https://aif.centre-mersenne.org/article/AIF_1993__43_5_1223_0.pdf}.
\bibitem{KP} K. Kurdyka and L. Paunescu, \emph{Hyperbolic polynomials and multiparameter real analytic perturbation theory}, arXiv:math/0602538. \url{https://arxiv.org/pdf/math/0602538}.
\bibitem{Lemmens} A. Lemmens, \emph{A local Jordan--Brouwer separation theorem}, arXiv:1810.13221, 2018. \url{https://arxiv.org/pdf/1810.13221}.
\bibitem{MST} M. Mazzucchelli, M. Salo, and L. Tzou, \emph{A general support theorem for analytic double fibration transforms}, arXiv:2306.05906. \url{https://users.jyu.fi/~salomi/pub/MazzucchelliSaloTzou2023.pdf}.
\bibitem{ML} K. W. Meng, M. H. Li, W. F. Yao, and X. Q. Yang, \emph{Lipschitz-like property relative to a set and the generalized Mordukhovich criterion}, arXiv:2009.10536. \url{https://arxiv.org/pdf/2009.10536}.
\bibitem{NS} M. Nedic and E. Saksman, \emph{On the structure of Nevanlinna measures}, arXiv:2201.01039. \url{https://arxiv.org/pdf/2201.01039}.
\bibitem{Ntalam} D. Ntalampekos, \emph{Monotone Sobolev functions in planar domains: level sets and smooth approximation}, arXiv:1911.06619. \url{https://arxiv.org/pdf/1911.06619}.
\bibitem{RW} R. T. Rockafellar and R. J.-B. Wets, \emph{Variational Analysis}, Springer, 1998. \url{https://sites.math.washington.edu/~rtr/papers/rtr169-VarAnalysis-RockWets.pdf}.
\bibitem{Schultka} K. Schultka, \emph{Microlocal analyticity of Feynman integrals}, dissertation, Humboldt-Universit\"at zu Berlin, 2019. \url{https://inspirehep.net/files/7007bb48bba430e220dafe390cae3c77}.
\bibitem{SW} A. Strohmaier and E. Witten, \emph{Analytic states in quantum field theory on curved spacetimes}, Annales Henri Poincar\'e \textbf{25} (2024), 4543--4590. \url{https://link.springer.com/article/10.1007/s00023-024-01419-0}.
\end{thebibliography}
\end{document}
